 % 本文件由 scripts/assemble.py 自动装配，勿手改（改 parts/ 后重跑）
% 第1章 狭义相对论与平直时空

% 第1章 Special Relativity and Flat Spacetime
\chapter{狭义相对论与平直时空}

% 1.1 Prelude
\section{引子}

广义相对论（general relativity，GR）是 Einstein 关于空间、时间和引力的理论。究其本质，这是一个非常简单的学科（比方说，与任何涉及量子力学的东西相比）。其核心思想再直截了当不过：自然界中的大多数力都用定义在时空上的场来表示（比如电磁场，或者亚核力所特有的短程场），而引力却内禀于时空自身之中。特别地，我们所体验到的“引力”，正是时空\emph{曲率}（curvature）的一种表现。

这样看来，我们的任务是清楚的。我们需要理解时空，需要理解曲率，还需要理解曲率如何变成引力。大致说来，本书前两章致力于探讨时空，第三章讨论曲率，第四章阐释曲率与引力之间的关系，然后才进入这个理论的各种应用。不过，让我们先放纵自己一下，对将要到来的内容做一番简短的预览，这或许会为我们旅程的头几步提供一点动力。

GR是关于引力的理论，所以不妨从回忆我们先前拥有的引力理论——Newton的理论——入手。它有两个基本要素：一个描述引力场如何受物质影响的方程，和一个描述物质如何响应这个场的方程。这些规则的Newton式传统表述用的是粒子之间的力：质量为$M$和$m$、由矢量$\mathbf{r}=r\,\mathbf{e}_{(r)}$相隔的两个物体之间的力，就是著名的平方反比定律，

\begin{equation*}
\mathbf{F} = -\frac{GMm}{r^2}\,\mathbf{e}_{(r)},
\tag{1.1}
\end{equation*}

这个力作用在质量为$m$的粒子上，按照Newton第二定律使它获得加速度，

\begin{equation*}
\mathbf{F} = m\mathbf{a}.
\tag{1.2}
\end{equation*}

等价地，我们也可以使用引力势$\Phi$的语言；势与质量密度$\rho$由Poisson方程相联系，

\begin{equation*}
\nabla^2\Phi = 4\pi G\rho,
\tag{1.3}
\end{equation*}

而加速度由势的梯度给出，

\begin{equation*}
\mathbf{a} = -\nabla\Phi.
\tag{1.4}
\end{equation*}

式 (1.1) 与式 (1.2)，或者式 (1.3) 与式 (1.4)，都足以定义Newton引力。而要定义GR，我们需要把它们中的每一组都替换成关于时空曲率的表述。

困难的部分是那个支配时空曲率对物质与能量之响应的方程。我们最终会以 Einstein 方程的形式找到想要的东西，

\begin{equation*}
R_{\mu\nu} - \frac{1}{2}R\,g_{\mu\nu} = 8\pi G\,T_{\mu\nu}.
\tag{1.5}
\end{equation*}

这个方程看上去比它本应有的样子更令人生畏，这多半得怪那些希腊下标。事实上，它只不过是$4\times4$矩阵之间的一个方程，下标标记的是每个矩阵的元素。左边的表达式量度时空的曲率，右边的则量度物质的能量和动量，所以正如所许诺的，这个方程把能量与曲率联系了起来。不过，对 Einstein 方程内部运作的详细理解，我们要留到以后再讲。

物质对时空曲率的响应则稍微容易把握一些：自由粒子沿着“可能的最短距离”的路径，也就是测地线（geodesic）运动。换句话说，粒子尽力沿直线运动，但在弯曲时空中可能根本不存在（我们在欧几里得几何中所熟悉的那种意义上的）直线，于是它们只能退而求其次。它们的参数化路径$x^\mu(\lambda)$服从测地线方程（geodesic equation）：

\begin{equation*}
\frac{d^2x^\mu}{d\lambda^2} + \Gamma^\mu_{\rho\sigma}\frac{dx^\rho}{d\lambda}\frac{dx^\sigma}{d\lambda} = 0.
\tag{1.6}
\end{equation*}

到了这一步，你对式 (1.6) 的理解不会比对式 (1.5) 多多少；但用不了多久，这一切都会豁然开朗。

我们后面将会讨论，测地线运动的普适性是GR极其深刻的一个特征。这种普适性正是我们如下论断的根源：引力其实并不是一种“力”，而是时空的一种性质。电场中的带电粒子会感受到一个加速度，使它偏离直线运动；与此相反，引力场中的粒子沿着的是世上最接近直线的东西。这样的粒子感受不到加速度；它们在自由下落。一旦我们更加熟悉GR的精神，下面这个想法就会显得完全合理：在空中飞行的球，比放在桌子上的球更真正地“未被加速”；放在桌上的那个球正在被推离它本来愿意所处的测地线（这就是为什么我们站在地球上时双脚会感受到力）。

我们描述时空曲率所依据的基本概念，将是度规张量（metric tensor），通常记作$g_{\mu\nu}$。度规通过表达对毕达哥拉斯定理$(\Delta l)^2=(\Delta x)^2+(\Delta y)^2$的偏离来编码一个空间的几何（其中$\Delta l$是定义在笛卡尔（Cartesian）网格上的两点之间的距离，两点的坐标间隔为$\Delta x$和$\Delta y$）。这个熟悉的公式只在传统的欧几里得几何中有效，那里隐含地假定了空间是平直的。一旦曲率出现，我们根深蒂固的几何观念就会开始失灵，而我们可以通过追踪毕达哥拉斯关系如何被改变来刻画曲率的大小。这一信息包含在度规张量之中。从度规出发，我们将推导出Riemann曲率张量——既用它来定义 Einstein 方程，也用它得到测地线方程。搭建这套数学装置，就是接下来几章的主题。

尽管为了定量地讨论曲率，不得不引入一定数量的形式语言，GR的核心观念（“引力就是时空的曲率”）却是相当简单的。那么，为什么GR至少在某些蒙昧的圈子里落得了困难、甚至玄奥的名声？因为 Einstein 理论的优雅真理被某些前相对论观念的层层累积所遮蔽；这些观念尽管非常有用，却必须先被抛弃，我们才能按GR的面貌领会这个世界。具体说来，我们生活在一个时空曲率非常小、而且粒子相对于光速大多运动得相当缓慢的世界里。因此，Galileo和Newton的力学对我们来说来得十分自然，尽管它只是那个更深的故事的一种近似。

所以，我们将着手学习这个更深的故事，办法是逐渐剥去那些有用却误导人的Newton直觉的层层外皮。第一步——也就是本章的主题——是探讨狭义相对论（special relativity，SR），即不存在引力（曲率）时的时空理论。希望这部分内容对读者来说大部分是复习，因为我们接下来会走得有些快。要点一方面在于重温SR究竟讲的是什么，另一方面在于引进张量和若干随后至关重要的相关概念，而不必同时承受曲率压在其他一切东西之上的额外复杂性。因此，本章我们将始终在平直时空中工作，而且只使用惯性（类似笛卡尔的）坐标。不用说，在你喜欢的任何坐标系里做SR都是可能的，但事实表明，为此引进必要的工具本身就会把我们带到通往弯曲空间的半路上，所以我们把这件事推迟一段时间。

% 1.2 Space and Time, Separately and Together
\section{空间与时间：分开来看与合起来看}

一种纯粹冷血的GR讲授方式会把第2章（流形）和第1章（狭义相对论与平直时空）的顺序颠倒过来。\emph{流形}（manifold）是用来描述时空的那类数学结构，而\emph{狭义相对论}则是一个调用某种特定时空（没有曲率、因而没有引力的时空）的模型。不过，既然你在读这本书，大概对狭义相对论（SR）至少有些熟悉，对流形却可能一无所知。所以我们的第一步是探访相对熟悉的SR领地，并利用这个机会引入一些对后续发展至关重要的概念和记号。

狭义相对论是关于时空结构的理论，而时空是粒子和场在其中演化的背景。SR充当了Newton力学的替代品，后者同样是一个关于时空结构的理论。无论哪种情形，我们都可以把这个基本结构与支配具体系统的各种动力学定律区分开来：Newton引力是置于Newton力学框架之内的一个动力学系统的例子，而麦克斯韦的电磁学则是在狭义相对论框架内运作的一个动力学系统。

\begin{figure}[H]
\centering
\includegraphics[width=0.72\textwidth]{fig_1.1.pdf}
\caption*{图 1.1\quad 在Newton时空中存在一种绝对的切片，把不同时刻的空间切成各不相同的拷贝。粒子世界线被约束只能随时间向前运动，但能以任意速度穿越空间；对于空间中不同地点的两个事件是否发生在同一时刻这个问题，存在人人一致的答案。}
\end{figure}

\textbf{时空}（spacetime）是一个四维集合，其元素由三个空间维度和一个时间维度来标记。（下一章我们会更严格地处理这些定义。）时空中的单独一个点称为一个\textbf{事件}（event）。粒子的路径是穿过时空的一条曲线，即一个参数化的一维事件集合，称为\textbf{世界线}（worldline）。这样的描述对SR和Newton力学同样适用。无论哪种情形，“时间”所受的对待显然与“空间”有些不同；具体地说，粒子总是沿时间向前行进，而在空间中却可以自由地来回移动。

然而，SR中粒子可以走的路径集合与Newton理论中的相比，有一个重要差别。在Newton力学中，时空有一个基本划分，分成轮廓分明的“固定时刻的全部空间”切片。\emph{同时性}（simultaneity）的观念——两个事件发生在同一时刻——是毫不含糊地定义了的。粒子的轨迹在时间上永远向前，但除此之外不受约束；特别是，两个这类粒子之间的相对速度没有上限。

在SR中，情况发生了戏剧性的改变：特别是，\emph{两个分离的事件“在同时”发生，这一观念没有良定义。}这并不是说时空完全没有什么结构。毋宁说，在每个事件处我们都可以定义一个\textbf{光锥}（light cone），它是穿过时空的那些路径的轨迹，这些路径可能被穿过该事件的光线所取。Newton力学中把时空唯一地切成由时间参数化的空间切片的那种绝对划分，被一条规则取代：物理粒子不能跑得比光快，因而其运动路径始终留在这些光锥之内。

\begin{figure}[H]
\centering
\includegraphics[width=0.72\textwidth]{fig_1.2.pdf}
\caption*{图 1.2\quad 在狭义相对论中，不存在"某一时刻的全部空间"这种绝对观念。取而代之的是一条规则：粒子永远以不大于光速的速率行进。于是我们可以在每个事件处定义光锥，它们局域地刻画了允许轨迹的集合。对于彼此的光锥之外的两个事件，哪个事件在时间上先发生，没有普适的观念。}
\end{figure}

SR中不存在偏好的时间切分，这正是时空观念在此语境中比在Newton力学中更为根本的关键所在。当然，我们可以在时空中选取具体的坐标系，而且一旦选定，在这个特定系统中谈论分离的事件取相同的时间坐标值就是有意义的；但还会存在其他可能的坐标系，它们通过把空间和时间相互“旋转”而与第一组相联系。这一现象是欧几里得几何中旋转的自然推广，我们现在就转向它。

考虑一块最寻常不过的二维平面。给这种平面上的点做标记，通常方便的做法是引入坐标，例如定义正交的$x$轴和$y$轴，按通常方式把每个点投影到这两个轴上。然而，显然关于这个平面的大多数有趣的几何事实都与我们坐标的选取无关；不存在任何偏好的方向。作为一个简单的例子，我们可以考虑两点之间的距离，它由下式给出：

\begin{equation*}
(\Delta s)^2 = (\Delta x)^2 + (\Delta y)^2.
\tag{1.7}
\end{equation*}

在一个不同的笛卡尔坐标系——由相对原来的轴旋转过的$x'$轴和$y'$轴定义——之中，距离的公式丝毫未变：

\begin{equation*}
(\Delta s)^2 = (\Delta x')^2 + (\Delta y')^2.
\tag{1.8}
\end{equation*}

因此我们说，距离在这样的坐标变更下不变。

\begin{figure}[H]
\centering
\includegraphics[width=0.72\textwidth]{fig_1.3.pdf}
\caption*{图 1.3\quad 二维欧几里得空间，画有两个不同的坐标系。"两点之间的距离"这类观念与所选的坐标系无关。}
\end{figure}

这就是为什么把平面看作一个内禀的二维空间是有用的，而不是把它看作两个被任意拼到一起的、从根本上不同的一维空间：虽然我们用两个不同的数来标记每个点，这些数并不是几何的本质，因为我们可以在保持距离不变的同时把轴相互旋转。在Newton物理学中，空间和时间却不是这样；把空间和时间相互旋转没有什么有用的含义。相反，“单一时刻的全部空间”这一观念具有独立于坐标的意义。

SR则是另一回事。让我们考虑时空上的坐标$(t,x,y,z)$，按如下方式建立。空间坐标$(x,y,z)$构成一个标准的笛卡尔系，例如可以把相交成直角的刚性杆焊接在一起来构造。这些杆必须自由运动、不受加速。时间坐标由一组时钟定义，这些时钟相对于空间坐标不运动。（既然这是思想实验，我们可以想象这些杆无限长，而且空间中每个点都有一只时钟。）这些时钟按下述意义同步。想象我们从空间中的点1沿直线以恒定速度$c$发送一束光到点2，然后立即回到1（速度为$-c$）。那么，光束到达点2时坐标钟上的时间（记作$t_2$），应当介于光束离开点1时坐标钟上的时间（$t_1$）与它回到同一只钟时的时间（$t_1'$）的正中间：

\begin{equation*}
t_2 = \frac{1}{2}(t_1' + t_1).
\tag{1.9}
\end{equation*}

如此构造出来的坐标系称为一个\textbf{惯性系}（inertial frame），或简称“惯性坐标”。这些坐标是空间中笛卡尔（正交归一）坐标向时空的自然推广。（如此小心翼翼地构造，是为了我们只在\emph{局域}上作比较；例如，绝不在同一时刻去比较两只相距遥远的钟。等到进入广义相对论，这种谨慎将变得更加必要，因为在那里将没有任何办法在整个时空上构造惯性坐标。）

\begin{figure}[H]
\centering
\includegraphics[width=0.72\textwidth]{fig_1.4.pdf}
\caption*{图 1.4\quad 在惯性坐标系中校准时钟。当一束光从点1射向点2再反弹回来时，如果时间$t_2$恰好位于$t_1$与$t_1'$的正中间，这些钟就是同步的。}
\end{figure}

我们可以用这套流程构造任意多个惯性系，它们与第一个的差别在于初始位置和时间的偏移、角度以及（恒定的）速度。在Newton式的世界里，新坐标$(t',x',y',z')$会有这样的特点：$t'=t+\text{常数}$，与空间坐标无关。也就是说，“两个事件同时发生，即在同一时刻发生”有一个绝对的含义。但在SR中这并不成立；一般而言，由$t=\text{常数}$定义的三维“空间”将不同于由$t'=\text{常数}$定义的那些。

不过，我们并没有完全坠入混乱。暂且不加任何动机，考虑一下我们将称之为两个事件之间的\textbf{时空间隔}（spacetime interval）的东西：

\begin{equation*}
\boxed{(\Delta s)^2 = -(c\Delta t)^2 + (\Delta x)^2 + (\Delta y)^2 + (\Delta z)^2.}
\tag{1.10}
\end{equation*}

（注意，即使对两个不重合的点，它也可以为正、为负或为零。）这里的$c$是空间与时间之间某个固定的换算因子，也就是说，一个固定的速度。作为经验事实，电磁波恰以这个速度$c$在真空中传播，因此我们把它称作“光速”。然而，重要的并不是光子碰巧以这个速率行进，而是存在这样一个$c$，使得\emph{时空间隔在惯性坐标的变更下不变}。换句话说，如果我们建立一个新的惯性系$(t',x',y',z')$，间隔将具有同样的形式：

\begin{equation*}
(\Delta s)^2 = -(c\Delta t')^2 + (\Delta x')^2 + (\Delta y')^2 + (\Delta z')^2.
\tag{1.11}
\end{equation*}

这正是为什么把SR看作四维时空的理论是有意义的，这个四维时空称为\textbf{Minkowski 空间}（Minkowski space）。（它是四维流形的一个特殊情形，流形我们稍后会详细处理。）我们将会看到，我们已经隐式定义的那些坐标变换，确实在某种意义上把空间和时间相互旋转。不存在“同时事件”的绝对观念；两件事是否发生在同一时刻，取决于所用的坐标。因此，把 Minkowski 空间分割成空间和时间，是我们为自己的目的而做的选择，并不是情形本身固有的东西。

与SR相联系的“佯谬”几乎全部源于对Newton观念的顽固固守：坚持存在唯一的时间坐标，坚持存在“单一时刻的全部空间”。改用时空来思考，而不是把空间与时间合在一起思考，这些佯谬往往就会消失。

让我们引入一些方便的记号。时空上的坐标将用带希腊字母上标指标的字母来表示，指标从0取到3，其中0一般表示时间坐标。于是，

\begin{equation*}
x^{\mu}:\quad
\begin{aligned}
x^0 &= ct \\
x^1 &= x \\
x^2 &= y \\
x^3 &= z.
\end{aligned}
\tag{1.12}
\end{equation*}

（不要开始把这些上标当成幂指数。）此外，为简单起见，我们将选取使下式成立的单位制：

\begin{equation*}
c = 1;
\tag{1.13}
\end{equation*}

因此在后面所有的公式里，我们都将省去因子$c$。经验告诉我们，$c$为$3\times10^8$米每秒；于是，我们是在1秒等于$3\times10^8$米的单位制下工作。有时分别指称$x^\mu$的空间分量和时间分量会很有用，所以我们将用拉丁字母上标单独表示空间分量：

\begin{equation*}
x^{i}:\quad
\begin{aligned}
x^1 &= x \\
x^2 &= y \\
x^3 &= z.
\end{aligned}
\tag{1.14}
\end{equation*}

把时空间隔写成更紧凑的形式也很方便。为此我们引入一个$4\times4$矩阵，即\textbf{度规}（metric），用两个下标书写：

\begin{equation*}
\eta_{\mu\nu} =
\begin{pmatrix}
-1 & 0 & 0 & 0 \\
0 & 1 & 0 & 0 \\
0 & 0 & 1 & 0 \\
0 & 0 & 0 & 1
\end{pmatrix}.
\tag{1.15}
\end{equation*}

（一些文献——尤其是场论教材——用相反的符号定义度规，所以要当心。）于是我们有漂亮的公式

\begin{equation*}
(\Delta s)^2 = \eta_{\mu\nu}\Delta x^\mu \Delta x^\nu.
\tag{1.16}
\end{equation*}

\begin{figure}[H]
\centering
\includegraphics[width=0.72\textwidth]{fig_1.5.pdf}
\caption*{图 1.5\quad 画在时空图上的一个光锥。图中标出了与原点类空分离、类光分离和类时分离的点。}
\end{figure}

这个公式引入了\textbf{求和约定}（summation convention）：既作为上标又作为下标出现的指标要对之求和。我们把这样的指标称为\textbf{哑指标}（dummy indices）；必须记住，它们对所有可能的值求和，而不取任何特定的值。（情形将永远如此：哑指标严格成对出现，一个“在上”，一个“在下”。这一点以后再谈。）因此，式 (1.16) 的内容与式 (1.10) 完全相同。

\textbf{时空图}（spacetime diagram）是一种极其有用的工具，那就让我们从这个观点来考察 Minkowski 空间。我们可以先把起初的$t$轴与$x$轴画成直角，并隐去$y$轴和$z$轴。（时空图上画出的“直角”并不必然意味着“在时空中正交”，尽管在这个例子里，$t$轴与$x$轴恰好确实正交。）考察以速率$c=1$行进的路径颇有启发，这些路径由$x=\pm t$给出。由以光速行进的直线与单个事件相连的全部点的集合，就是\textbf{光锥}；因为如果想象再纳入一个空间维度，那两条对角线就会补全成一个锥面。光锥天然分为未来光锥与过去光锥；一点$p$的未来与过去光锥内部所有点的集合，称为与$p$\textbf{类时分离}（timelike separated）；光锥之外的点与$p$\textbf{类空分离}（spacelike separated）；而恰好位于锥面上的点，则与$p$\textbf{类光分离}（lightlike or null separated）。回到式 (1.10)，我们看到类时分离的点之间的间隔为负，类空分离的点之间的间隔为正，类光分离的点之间的间隔为零。（间隔定义为$(\Delta s)^2$本身，而不是这个量的平方根。）

间隔对类时曲线（比光慢的粒子实际上将沿之运动）为负，这件事有些恼人，所以我们定义\textbf{固有时}（proper time）$\tau$满足

\begin{equation*}
(\Delta\tau)^2 = -(\Delta s)^2 = -\eta_{\mu\nu}\Delta x^\mu \Delta x^\nu.
\tag{1.17}
\end{equation*}

时空间隔的一个关键性质是：\emph{两个事件之间的固有时，量度的是沿连接这两事件的直线路径运动的观者所经历的时间。}这在非常特殊的情形下最容易看清：两个事件具有相同的空间坐标，只在时间上分离；这对应于在两事件之间旅行的观者相对所用坐标系静止。于是$(\Delta\tau)^2=-\eta_{\mu\nu}\Delta x^\mu\Delta x^\nu=(\Delta t)^2$，故$\Delta\tau=\Delta t$；当然，我们本来就把$t$定义为位于固定空间位置上的钟所量度的时间。但时空间隔在惯性系的变更下不变；因此，两个固定事件之间的固有时 (1.17) 在观者运动着的惯性系中求出的值，与在观者静止的惯性系中求出的值相同。

一个关键的事实是：对更一般的轨迹，固有时与坐标时并不相同（尽管固有时总是沿轨迹运动的观者所携带的钟量度的时间）。考虑连接事件$A$和$C$的两条轨迹：一条是直线，经过标记为$B$的中点；另一条由一位观者走完，他以恒定速度$v=dx/dt$从$A$离开行进到一点$B'$，再以恒定速度$-v$返回，与前者相交于事件$C$。

\begin{figure}[H]
\centering
\includegraphics[width=0.72\textwidth]{fig_1.6.pdf}
\caption*{图 1.6\quad 双生子佯谬。沿穿过时空的直线路径$ABC$旅行的行者，将比沿非直线路径$AB'C$旅行的人衰老得更多。既然固有时是穿过时空走过的距离的量度，这应当不足为怪。（唯一可能令人惊讶的是，直线路径竟然是固有时\emph{最长}的那一条；这可以追溯到度规类时分量的负号。）}
\end{figure}

选取惯性坐标，使直的轨迹描述一个静止不动的粒子，事件$A$位于坐标$(t,x)=(0,0)$，$C$位于$(\Delta t,0)$。两条路径于是在时空中构成一个等腰三角形；$B$的坐标为$(\frac{1}{2}\Delta t,0)$，$B'$的坐标为$(\frac{1}{2}\Delta t,\Delta x)$，其中$\Delta x=\frac{1}{2}v\Delta t$。显然有$\Delta\tau_{AB}=\frac{1}{2}\Delta t$，但是

\begin{equation*}
\begin{aligned}
\Delta\tau_{AB'} &= \sqrt{\left(\tfrac{1}{2}\Delta t\right)^2 - (\Delta x)^2} \\
&= \tfrac{1}{2}\sqrt{1-v^2}\,\Delta t.
\end{aligned}
\tag{1.18}
\end{equation*}

应当显然有$\Delta\tau_{BC}=\Delta\tau_{AB}$和$\Delta\tau_{B'C}=\Delta\tau_{AB'}$。于是，沿直线从事件$A$旅行到$C$的观者经历的时间流逝为$\Delta\tau_{ABC}=\Delta t$，而走出去又返回的那位经历的却是

\begin{equation*}
\Delta\tau_{AB'C} = \sqrt{1-v^2}\,\Delta t < \Delta t.
\tag{1.19}
\end{equation*}

尽管两位观者在时空中的同一点出发、又在同一点结束，他们衰老的程度却不一样。这就是著名的“双生子佯谬”（twin paradox）——形形色色的误解和蹩脚解释的不幸上演之处。真相很直接：时空中的非直线路径与直线路径有不同的间隔，正如空间中的非直线路径与直线路径有不同的长度。当然，这件事并不像听起来那么平凡；深刻的洞见在于“沿世界线流逝的时间”以何种方式与穿过时空历经的间隔相联系。在Newton式的世界里，坐标$t$代表贯穿全部时空的普适的时间之流；在相对论中，$t$只是一个方便的坐标，时间的流逝取决于你沿哪条路径旅行。一个重要的区别是：非直线路径的固有时\emph{更短}。在空间中，两点之间最短的距离是直线；在时空中，两事件之间最长的固有时属于直的轨迹。

并非所有轨迹都齐整到可以由直线段拼成。在更一般的情形下，引入无穷小间隔，即\textbf{线元}（line element），会很有用：

\begin{equation*}
\boxed{ds^2 = \eta_{\mu\nu}dx^\mu dx^\nu,}
\tag{1.20}
\end{equation*}

其中$dx^\mu$是无穷小的坐标位移。（我们这里的处理相当不正式，不过以后会补正。）从这个定义出发，人们很想取平方根并沿一条路径积分，以得到有限的间隔，但$\int\sqrt{\eta_{\mu\nu}dx^\mu dx^\nu}$究竟应当是什么意思，多少有些不清楚。换一种做法，我们把穿过时空的一条路径看作参数化曲线$x^\mu(\lambda)$。注意，与Newton力学的传统做法不同，参量$\lambda$不一定等同于时间坐标。于是我们可以计算导数$dx^\mu/d\lambda$，并把沿类空曲线（无穷小间隔为类空的曲线）的路径长度写作

\begin{equation*}
\Delta s = \int \sqrt{\eta_{\mu\nu}\frac{dx^\mu}{d\lambda}\frac{dx^\nu}{d\lambda}}\,d\lambda,
\tag{1.21}
\end{equation*}

其中积分沿整条路径进行。对类时路径我们用固有时

\begin{equation*}
\Delta\tau = \int \sqrt{-\eta_{\mu\nu}\frac{dx^\mu}{d\lambda}\frac{dx^\nu}{d\lambda}}\,d\lambda,
\tag{1.22}
\end{equation*}

它将为正。（对类光路径，间隔干脆就是零。）当然，我们可以考虑在一些地方类时、在另一些地方类空的路径，但幸运的是很少有必要，因为物理粒子的路径从不改变其特征（有质量粒子沿类时路径运动，无质量粒子沿类光路径运动）。再说一遍，$\Delta\tau$确实就是沿轨迹运动的观者所量度的时间。

狭义相对论中的\emph{加速度}（acceleration）这一概念名声不佳，但毫无道理。在建立惯性坐标时，我们\emph{当然}小心翼翼地确保了在此类坐标中静止的粒子不受加速。然而，一旦建立了这样的坐标，我们就可以随意考虑物理粒子的任何轨迹，无论加速与否。特别是，“SR处理不了加速轨迹、必须请出广义相对论”的传言毫无事实依据。广义相对论是在有引力之时、当时空变得弯曲之际才变得相关。平直时空中的任何过程都在狭义相对论的框架内得到描述；特别地，像式 (1.22) 这样的表达式是完全普遍的。

% 1.3 Lorentz Transformations
\section{洛伦兹变换}

现在，我们可以在比之前稍抽象一点的层次上考虑时空中的坐标变换。我们感兴趣的是如何把经由上述流程构造的各种惯性系联系起来，并对此给出一个形式化的描述；也就是说，（考虑那些）使间隔 (1.16) 保持不变的坐标系。其中简单的一类是\textbf{平移}（translations），它们只是让坐标平移（在空间或时间上）：

\begin{equation*}
x^\mu \rightarrow x^{\mu'} = \delta^{\mu'}_{\mu}\,(x^\mu + a^\mu),
\tag{1.23}
\end{equation*}

其中$a^\mu$是由四个固定数构成的一组数，而$\delta^{\mu'}_{\mu}$是传统Kronecker delta符号的四维版本：

\begin{equation*}
\delta^{\mu'}_{\mu} =
\begin{cases}
1 & \text{当 } \mu' = \mu, \\
0 & \text{当 } \mu' \neq \mu.
\end{cases}
\tag{1.24}
\end{equation*}

注意，我们把撇号加在指标上，而不是加在$x$上。这样做的理由，等我们开始同矢量和张量打交道之后会变得更加清楚；这种记号提醒我们：几何对象是同一个，只是它的分量是相对于另一个坐标系分解的。平移保持差值$\Delta x^\mu$不变，所以间隔不变并不令人惊讶。其他相关的变换还包括空间的\textbf{旋转}（rotations），以及以一个恒定速度矢量做的偏移，即\textbf{推动}（boosts）；这些都是线性变换，可以用一个（不依赖时空的）矩阵去乘$x^\mu$来描述：

\begin{equation*}
x^{\mu'} = \Lambda^{\mu'}{}_{\nu}x^\nu,
\tag{1.25}
\end{equation*}

或者，用更传统的矩阵记号，

\begin{equation*}
x' = \Lambda x.
\tag{1.26}
\end{equation*}

（我们一般会使用指标，而不是矩阵记号，但此刻我们有意把这里的讨论与通常用矩阵描述的某些其他熟悉概念联系起来。）这些变换并不使差值$\Delta x^\mu$保持不变，而是把它们再乘上矩阵$\Lambda$。什么样的矩阵能使间隔不变？坚持用矩阵记号的话，我们想要的是

\begin{equation*}
\begin{aligned}
(\Delta s)^2 &= (\Delta x)^{\mathrm{T}}\eta(\Delta x) = (\Delta x')^{\mathrm{T}}\eta(\Delta x') \\
&= (\Delta x)^{\mathrm{T}}\Lambda^{\mathrm{T}}\eta\Lambda(\Delta x),
\end{aligned}
\tag{1.27}
\end{equation*}

因此

\begin{equation*}
\eta = \Lambda^{\mathrm{T}}\eta\Lambda,
\tag{1.28}
\end{equation*}

% 本文件至此为止：书页12末句止于式 (1.28) 后的逗号，正文延续到书页13，由下一块（ch1_b）接续。
或
\begin{equation*}
\eta_{\rho\sigma}=\Lambda^{\mu'}{}_{\rho}\,\eta_{\mu'\nu'}\,\Lambda^{\nu'}{}_{\sigma}
=\Lambda^{\mu'}{}_{\rho}\Lambda^{\nu'}{}_{\sigma}\eta_{\mu'\nu'}\,.
\tag{1.29}
\end{equation*}

（在矩阵记号中次序是重要的，而在指标记号中则无关紧要。）我们想要找出矩阵$\Lambda^{\mu'}{}_{\nu}$，使得矩阵$\eta_{\mu'\nu'}$的分量与$\eta_{\rho\sigma}$的分量相同；这就是间隔在这些变换下不变的含义。

满足式(1.28)的那些矩阵称为\textbf{洛伦兹变换}（Lorentz transformations）；它们的全体在矩阵乘法下构成一个群，称为\textbf{洛伦兹群}（Lorentz group）。这个群与SO(3)——三维空间中的转动群——之间有着密切的类比。转动群可以看成由满足$R^{\mathrm{T}}\mathbf{1}R=\mathbf{1}$的$3\times3$矩阵$R$组成，其中$\mathbf{1}$是$3\times3$单位矩阵。这样的矩阵称为\emph{正交}（orthogonal）矩阵，$3\times3$的这种矩阵构成群O(3)。其中不仅有转动，还包括空间坐标轴取向的反转（\textbf{宇称变换}（parity transformations））。有时我们还想排除宇称变换，于是进一步要求矩阵的行列式为1（$|R|=1$）；这样的矩阵称为\emph{特殊}（special）矩阵，所得的群就是SO(3)。如果把它写成

\begin{equation*}
\mathbf{1}=R^{\mathrm{T}}\mathbf{1}R\,,
\tag{1.30}
\end{equation*}

正交条件看上去就更像式(1.28)了。

于是，转动群O(3)与洛伦兹群的差别就在于：把$\mathbf{1}$——所有元素都等于$+1$的$3\times3$对角矩阵——换成$\eta$，一个对角元一个为$-1$、其余为$+1$的$4\times4$对角矩阵。因此洛伦兹群常被称作O(3,1)。它不仅包含推动和转动，还包含时间方向的分立反转以及宇称变换。和前面一样，我们可以要求$|\Lambda|=1$，得到“固有洛伦兹群”（proper Lorentz group）SO(3,1)。然而这并没有给我们真正想要的东西，即连续洛伦兹变换的集合（那些与恒等变换光滑连通的变换），因为时间反转与宇称反转的组合会具有单位行列式。从式(1.29)的$(\rho,\sigma)=(0,0)$分量出发，我们容易证明$|\Lambda^{0'}{}_{0}|\geq1$，其中负值对应时间反转。因此我们最终可以再要求$\Lambda^{0'}{}_{0}\geq1$（加上$|\Lambda|=1$），得到“固有正向”（proper orthochronous）或“受限”（restricted）洛伦兹群。有时它被记作类似$SO(3,1)\uparrow$的符号，但通常我们不会费心去明确作出这种区分。注意，$3\times3$单位矩阵其实就是普通平直空间的度规。所有本征值都为正的这种度规称为\textbf{欧几里得型}（Euclidean）度规，而像式(1.15)那样只带一个负号的度规则称为\textbf{洛伦兹型}（Lorentzian）度规。

写下简单洛伦兹变换的显式表达式是直截了当的。$x$-$y$平面中一个我们熟悉的转动是

\begin{equation*}
\Lambda^{\mu'}{}_{\nu}=
\begin{pmatrix}
1 & 0 & 0 & 0\\
0 & \cos\theta & \sin\theta & 0\\
0 & -\sin\theta & \cos\theta & 0\\
0 & 0 & 0 & 1
\end{pmatrix}.
\tag{1.31}
\end{equation*}

转动角$\theta$是周期为$2\pi$的周期变量。推动则可以想成“空间方向与时间方向之间的转动”。沿$x$方向的推动就是一个例子：

\begin{equation*}
\Lambda^{\mu'}{}_{\nu}=
\begin{pmatrix}
\cosh\phi & -\sinh\phi & 0 & 0\\
-\sinh\phi & \cosh\phi & 0 & 0\\
0 & 0 & 1 & 0\\
0 & 0 & 0 & 1
\end{pmatrix}.
\tag{1.32}
\end{equation*}

推动参量$\phi$与转动角不同，其定义范围是从$-\infty$到$\infty$。把各个变换相乘就可以得到一般的变换；这个六参数矩阵（三个推动、三个转动）的显式表达式并不好看，也不值得费神写下来。一般说来洛伦兹变换不可对易，所以洛伦兹群是非阿贝尔群（nonabelian）。平移与洛伦兹变换合在一起构成一个十参数非阿贝尔群，即\textbf{庞加莱群}（Poincaré group）。

大家得知推动对应于通过换到一个以恒定速度运动的参考系来改变坐标时，不应感到惊讶，但我们不妨更明白地看一下。（不要把“推动”（boosting）与“加速”（accelerating）混淆。变到一个不同参考系的推动与加速一个物体这两者之间的差别，就如同转到不同坐标系与让物体自旋这两者之间的差别。）对于由式(1.32)给出的变换，变换后的坐标$t'$和$x'$将为

\begin{align*}
t'&=t\cosh\phi-x\sinh\phi\\
x'&=-t\sinh\phi+x\cosh\phi\,.
\tag{1.33}
\end{align*}

由此我们看到，由$x'=0$定义的那个点在运动；它具有速度

\begin{equation*}
v=\frac{x}{t}=\frac{\sinh\phi}{\cosh\phi}=\tanh\phi\,.
\tag{1.34}
\end{equation*}

为了换成更平实的记号，我们可以代入$\phi=\tanh^{-1}v$，得到

\begin{align*}
t'&=\gamma(t-vx)\\
x'&=\gamma(x-vt)\,,
\tag{1.35}
\end{align*}

其中$\gamma=1/\sqrt{1-v^{2}}$。的确如此，我们这条抽象路线重新得到了洛伦兹变换的惯用表达式。应用这些公式就得到时间膨胀（time dilation）、长度收缩（length contraction）等等。

在时空图（spacetime diagram）的语境中考察洛伦兹变换是很有启发性的。根据式(1.33)，在$x$-$t$平面内的推动下，$x'$轴（$t'=0$）由$t=x\tanh\phi$给出，而$t'$轴（$x'=0$）则由$t=x/\tanh\phi$给出。于是我们看到，空间轴和时间轴彼此向对方转过去，只不过它们像剪刀一样斜着合拢，而不再保持传统欧几里得意义下的正交。（我们将会看到，这些轴事实上在洛伦兹意义下确实仍保持正交；这正是度规在推动下保持不变的含义。）

\begin{figure}[H]
\centering
\includegraphics[width=0.72\textwidth]{fig_1.7.pdf}
\caption*{图 1.7\quad 洛伦兹变换把$\{t',x'\}$坐标与$\{t,x\}$坐标联系起来。注意光锥不变。}
\end{figure}

这应该毫不奇怪，因为假如时空的行为就像四维版本的空间，那世界将会是一个非常不同的地方。我们相当鲜明地看到了这种情形与Newton世界之间的区别；在狭义相对论中，我们不可能（以不依赖于坐标的方式）说出一个与点$p$类空间隔的点究竟是在$p$的未来、$p$的过去，还是“与$p$同时”。

还要注意，由$x'=\pm t'$定义的路径与由$x=\pm t$定义的路径完全相同；这些轨迹在沿$x$轴的推动下保持不变。当然，我们知道光正是以这个速度传播的；于是我们便发现，光速在任何惯性系中都是相同的。

\section{矢量}

要更细致地探究 Minkowski 空间的结构，就需要引入矢量与张量的概念。我们先从矢量讲起，大家应该对它已经熟悉。当然，时空中的矢量是四维的，常被称为\textbf{四维矢量}（four-vectors）。这会带来相当可观的差别——例如，两个四维矢量之间根本不存在叉积这种东西。

除了维数这个简单事实之外，最需要强调的一点是：每个矢量都定域在时空中的一个给定点上。你可能习惯于把矢量看成从空间中一点拉伸到另一点的东西，甚至看成可以随意从一个点滑动到另一个点的“自由”矢量。在平直空间的语境之外，这些都不是有用的概念；一旦引入曲率，我们就失去了从一个点向另一个点画出特殊曲线的能力，也失去了把矢量唯一地搬动绕行流形的能力。毋宁说，对时空中的每个点$p$，我们都把定域于该点的所有可能矢量的集合与它联系起来；这个集合称为$p$处的\textbf{切空间}（tangent space），记作$T_p$。这个名称的灵感来自这样的图像：在一个简单的弯曲二维空间中，附着于某点的那些矢量全体构成一个在该点与此空间相切的平面。（这幅图像依赖于把流形和切空间嵌入一个更高维的外部空间，而一般来说我们既没有也不需要这种嵌入。）抛开这个灵感不谈，重要的是把这些矢量看成定域于单独一点，而不是从一点拉伸到另一点（尽管这不会妨碍我们在时空图上把它们画成箭头）。

在第2章中，我们将把每个点处的切空间与我们可以从时空本身构造出来的东西联系起来。眼下，只需把$T_p$想成时空每一点上的一个抽象矢量空间。所谓\textbf{（实）矢量空间}（real vector space），是指由一些对象（矢量）组成的集合，它们可以线性地彼此相加并乘以实数。于是，对任意两个矢量$V$、$W$和实数$a$、$b$，我们有

\begin{equation*}
(a+b)(V+W)=aV+bV+aW+bW\,.
\tag{1.36}
\end{equation*}

每个矢量空间都有一个原点，即在矢量加法下充当恒等元的零矢量。在许多矢量空间中还有别的运算，比如取内积（点积），但这是矢量空间这一初等概念之外的额外结构。

矢量是一个定义得非常明确的几何对象，\textbf{矢量场}（vector field）也一样，它定义为在时空每一点上恰好有一个矢量的矢量的集合。[把一个$n$维流形$M$的所有切空间汇集起来，可以组装成一个$2n$维流形，称为\textbf{切丛}（tangent bundle），$T(M)$。它是“纤维丛”（fiber bundle）的一个具体例子，纤维丛还带有某种额外的数学结构；就目前的目的我们不需要那些细节。]尽管如此，把矢量相对于某组基矢分解成分量往往是有用的。\textbf{基}（basis）是任何一组既张满矢量空间（任何矢量都是基矢的线性组合）又线性无关（基中没有任何一个矢量是其余基矢的线性组合）的矢量。对任何给定的矢量空间，可供选择的基有无穷多种，但每个基都由相同数目的矢量组成，这个数目称为该空间的\textbf{维数}（dimension）。（对于 Minkowski 空间中一点相应的切空间，维数当然是四。）

\begin{figure}[H]
\centering
\includegraphics[width=0.72\textwidth]{fig_1.8.pdf}
\caption*{图 1.8\quad 切空间$T_p$的示意性图示——点$p$处所有矢量构成的空间。}
\end{figure}

设想在每个切空间中都立起一组四个矢量$\hat{e}_{(\mu)}$作为基，$\mu$和通常一样取遍$\{0,1,2,3\}$。实际上，不妨说每个基都是“适配于坐标$x^\mu$”的——也就是说，基矢$\hat{e}_{(1)}$就是我们通常认为指向$x$轴方向的那个矢量。我们完全没有必要选择适配于任何坐标系的基，尽管这样做往往很方便。（和前面一样，我们在这里本可以更精确一些，但后面我们会以令人难以忍受的精确程度重新讨论这个问题，所以现在稍微马虎一点是可以原谅的。）于是，任何抽象矢量$A$都可以写成基矢的线性组合：

\begin{equation*}
A=A^{\mu}\hat{e}_{(\mu)}\,.
\tag{1.37}
\end{equation*}

系数$A^{\mu}$就是矢量$A$的\textbf{分量}（components）。我们多半会把基完全忘掉，而有点随意地把“矢量$A^{\mu}$”挂在嘴边，但要记住这只是一种简写。真正的矢量是抽象的几何实体，而分量只是矢量在某个方便的基下基矢的系数。（由于我们通常会不写出显式基矢，指标一般就用来标记矢量和张量的分量。这就是基矢的指标要加括号的原因——提醒我们这是矢量的集合，而不是单个矢量的分量。）

时空中的一个矢量的标准例子是曲线的切矢量。穿过时空的一条参数化曲线或路径由坐标作为参量的函数给出，例如$x^{\mu}(\lambda)$。切矢量$V(\lambda)$的分量为

\begin{equation*}
V^{\mu}=\frac{\mathrm{d}x^{\mu}}{\mathrm{d}\lambda}\,.
\tag{1.38}
\end{equation*}

完整的矢量是$V=V^{\mu}\hat{e}_{(\mu)}$。在洛伦兹变换下，坐标$x^{\mu}$按照式(1.25)变换，而参量化$\lambda$不变；因此我们可以推出，切矢量的分量必须按如下方式变化

\begin{equation*}
\boxed{\,V^{\mu}\rightarrow V^{\mu'}=\Lambda^{\mu'}{}_{\nu}V^{\nu}\,.}
\tag{1.39}
\end{equation*}

然而，矢量$V$本身（与其在某个坐标系中的分量相对而言）在洛伦兹变换下是不变的。我们可以利用这个事实来导出基矢的变换性质。把变换后坐标系中的那组基矢记为$\hat{e}_{(\nu')}$。由于矢量不变，我们有

\begin{equation*}
V=V^{\mu}\hat{e}_{(\mu)}=V^{\nu'}\hat{e}_{(\nu')}
=\Lambda^{\nu'}{}_{\mu}V^{\mu}\hat{e}_{(\nu')}\,.
\tag{1.40}
\end{equation*}

但无论分量$V^{\mu}$的数值是什么，这个关系都必须成立。因此我们可以说

\begin{equation*}
\hat{e}_{(\mu)}=\Lambda^{\nu'}{}_{\mu}\hat{e}_{(\nu')}\,.
\tag{1.41}
\end{equation*}

要把新基$\hat{e}_{(\nu')}$用旧基$\hat{e}_{(\mu)}$表示出来，我们应该乘以洛伦兹变换$\Lambda^{\nu'}{}_{\mu}$的逆。但是，从无撇坐标变到带撇坐标的洛伦兹变换的逆也是一个洛伦兹变换，只不过这次是从带撇系变到无撇系。因此我们将引入一种略显精妙的记号：对这两个矩阵用同一个符号，只是把带撇指标与无撇指标互换。也就是说，由$\Lambda^{\mu'}{}_{\nu}$规定的洛伦兹变换，其逆变换写作$\Lambda^{\rho}{}_{\sigma'}$。从运算上讲，这意味着

\begin{equation*}
\Lambda^{\mu}{}_{\nu'}\Lambda^{\nu'}{}_{\rho}=\delta^{\mu}{}_{\rho}\,,\qquad
\Lambda^{\sigma'}{}_{\lambda}\Lambda^{\lambda}{}_{\tau'}=\delta^{\sigma'}{}_{\tau'}\,.
\tag{1.42}
\end{equation*}

于是由式(1.41)我们得到基矢的变换规则：

\begin{equation*}
\hat{e}_{(\nu')}=\Lambda^{\mu}{}_{\nu'}\hat{e}_{(\mu)}\,.
\tag{1.43}
\end{equation*}

因此，基矢的集合按坐标或矢量分量的洛伦兹变换之逆来变换。

让我们稍停片刻，把这一切消化一下。我们引入了用上指标标记的坐标，它们在洛伦兹变换下按某种方式变换。接着我们考虑了也用上指标写出的矢量分量，这是有意义的，因为它们以与坐标函数相同的方式变换。（在固定的坐标系中，四个坐标$x^{\mu}$中的每一个都可以看成时空上的函数，矢量场的四个分量中的每一个也是如此。）与坐标系相联系的基矢则通过逆矩阵变换，并用下指标标记。这种记号保证了由分量与基矢求和构造出来的不变对象在变换下保持不变，正如我们所愿。如果我说，对于可能有多个指标的张量，情形也将继续如此，大概不算把后面的内容剧透太多。

\section{对偶矢量（1-形式）}

一旦建立了一个矢量空间，我们就可以定义另一个与之相联系的矢量空间（维数相同），称为\textbf{对偶矢量空间}（dual vector space）。对偶空间通常用星号标记，于是切空间$T_p$的对偶空间称为\textbf{余切空间}（cotangent space），记作$T_p^{*}$。对偶空间是从原来的矢量空间到实数的所有线性映射构成的空间；用数学行话说，若$\omega\in T_p^{*}$是一个对偶矢量（dual vector），则它作为一个映射来作用，满足

\begin{equation*}
\omega(aV+bW)=a\omega(V)+b\omega(W)\in\mathbf{R}\,,
\tag{1.44}
\end{equation*}

其中$V$、$W$是矢量，$a$、$b$是实数。这些映射的妙处在于它们本身也构成一个矢量空间；于是，若$\omega$和$\eta$是对偶矢量，我们有

\begin{equation*}
(a\omega+b\eta)(V)=a\omega(V)+b\eta(V)\,.
\tag{1.45}
\end{equation*}

为了让这个构造更具体一点，我们可以引入一组基对偶矢量$\hat{\theta}^{(\nu)}$，办法是要求

\begin{equation*}
\hat{\theta}^{(\nu)}(\hat{e}_{(\mu)})=\delta^{\nu}{}_{\mu}\,.
\tag{1.46}
\end{equation*}

于是每个对偶矢量都可以用它的分量写出，我们用下指标来标记分量：

\begin{equation*}
\omega=\omega_{\mu}\hat{\theta}^{(\mu)}\,.
\tag{1.47}
\end{equation*}

通常，与矢量的情形完全类似，我们会径直写$\omega_{\mu}$，用它代表整个对偶矢量。实际上，你有时会看到$T_p$的元素（就是我们所说的矢量）被称为\textbf{逆变矢量}（contravariant vectors），而$T_p^{*}$的元素（就是我们所说的对偶矢量）被称为\textbf{协变矢量}（covariant vectors），不过在如今这个时代，这些术语听起来有点过时了。只要你把普通矢量称作带上指标的矢量、把对偶矢量称作带下指标的矢量，应该没有人会不高兴。对偶矢量的另一个名字是\textbf{1-形式}（one-forms），这个多少有些神秘的称呼在第2章中会变得更清楚。

分量记号给出了一种书写对偶矢量作用于矢量的简单方式：

\begin{align*}
\omega(V)&=\omega_{\mu}\hat{\theta}^{(\mu)}(V^{\nu}\hat{e}_{(\nu)})\\
&=\omega_{\mu}V^{\nu}\hat{\theta}^{(\mu)}(\hat{e}_{(\nu)})\\
&=\omega_{\mu}V^{\nu}\delta^{\mu}{}_{\nu}\\
&=\omega_{\mu}V^{\mu}\in\mathbf{R}\,.
\tag{1.48}
\end{align*}

这正是我们很少需要把基矢和对偶矢量显式写出来的原因；分量包办了一切工作。式(1.48)的形式还暗示，通过定义

\begin{equation*}
V(\omega)\equiv\omega(V)=\omega_{\mu}V^{\mu}\,.
\tag{1.49}
\end{equation*}

我们可以把矢量看成对偶矢量上的线性映射。因此，对偶矢量空间的对偶空间就是原来的矢量空间本身。

当然，在时空中我们感兴趣的不是一个单独的矢量空间，而是矢量场和对偶矢量场。[把$M$上所有余切空间汇合起来，就得到\textbf{余切丛}（cotangent bundle），$T^{*}(M)$。]在那种情形下，对偶矢量场作用在矢量场上得到的不是一个单独的数，而是时空上的一个\textbf{标量}（scalar）（函数）。标量是不带指标的量，它在洛伦兹变换下不变；它是从时空到实数的、不依赖于坐标的映射。

我们可以用前面处理矢量时用过的同样论证（几何对象独立于坐标，即使它们的分量并非如此）来导出对偶矢量的变换性质。答案是：对于分量，

\begin{equation*}
\boxed{\,\omega_{\mu'}=\Lambda^{\nu}{}_{\mu'}\omega_{\nu}\,,}
\tag{1.50}
\end{equation*}

而对于基对偶矢量，

\begin{equation*}
\hat{\theta}^{(\rho')}=\Lambda^{\rho'}{}_{\sigma}\hat{\theta}^{(\sigma)}\,.
\tag{1.51}
\end{equation*}

这正是从指标位置所能预期的；对偶矢量的分量按矢量分量的逆变换来变换。注意这就保证了标量(1.48)在洛伦兹变换下不变，正如它本应如此。

在时空中，对偶矢量最简单的例子是标量函数的\textbf{梯度}（gradient），即对时空坐标的偏导数的集合，我们用小写的d来记它：

\begin{equation*}
\mathrm{d}\phi=\frac{\partial\phi}{\partial x^{\mu}}\hat{\theta}^{(\mu)}\,.
\tag{1.52}
\end{equation*}

用来变换偏导数的常规链式法则，在这种情形下恰好就是对偶矢量分量的变换法则：

\begin{align*}
\frac{\partial\phi}{\partial x^{\mu'}}&=\frac{\partial x^{\mu}}{\partial x^{\mu'}}\frac{\partial\phi}{\partial x^{\mu}}\\
&=\Lambda^{\mu}{}_{\mu'}\frac{\partial\phi}{\partial x^{\mu}}\,,
\tag{1.53}
\end{align*}

其中我们用了式(1.25)把洛伦兹变换与坐标联系起来。梯度是对偶矢量这一事实，导致了偏导数的如下简写记号：

\begin{equation*}
\frac{\partial\phi}{\partial x^{\mu}}=\partial_{\mu}\phi=\phi_{,\mu}\,.
\tag{1.54}
\end{equation*}

于是，$x^{\mu}$带上指标，但当它出现在导数的分母中时，它意味着所得的对象带下指标。在本书中我们一般使用$\partial_{\mu}$而不用逗号记号。注意，对于前面给出的那个矢量例子——曲线的切矢量，梯度确实以一种自然的方式作用。其结果是函数沿曲线的普通导数：

\begin{equation*}
\partial_{\mu}\phi\,\frac{\partial x^{\mu}}{\partial\lambda}=\frac{\mathrm{d}\phi}{\mathrm{d}\lambda}\,.
\tag{1.55}
\end{equation*}

\section{张量}

矢量与对偶矢量的一个直接推广就是\textbf{张量}（tensor）的概念。正如对偶矢量是从矢量到$\mathbf{R}$的线性映射，一个型（或阶）$(k,l)$的张量$T$是从一组对偶矢量和矢量到$\mathbf{R}$的多重线性映射（multilinear map）：

\begin{equation*}
T:\;T_p^{*}\times\underset{(k\text{ 次})}{\cdots}\times T_p^{*}
\times T_p\times\underset{(l\text{ 次})}{\cdots}\times T_p\rightarrow\mathbf{R}\,.
\tag{1.56}
\end{equation*}

这里，“$\times$”表示笛卡尔积（Cartesian product），例如$T_p\times T_p$就是有序矢量对构成的空间。多重线性意味着张量作用在它的每一个宗量上都是线性的；例如，对于型$(1,1)$的张量，我们有

\begin{align*}
T(a\omega+b\eta,\,cV+dW)={}&acT(\omega,V)\\
&+adT(\omega,W)+bcT(\eta,V)+bdT(\eta,W)\,.
\tag{1.57}
\end{align*}

从这个观点来看，标量是型$(0,0)$的张量，矢量是型$(1,0)$的张量，而对偶矢量是型$(0,1)$的张量。

固定型$(k,l)$的所有张量构成一个矢量空间；它们可以彼此相加并乘以实数。为了给这个空间构造基，我们需要定义一种新的运算，称为\textbf{张量积}（tensor product），记作$\otimes$。若$T$是$(k,l)$张量，$S$是$(m,n)$张量，我们如下定义$(k+m,\,l+n)$张量$T\otimes S$：

\begin{align*}
&T\otimes S(\omega^{(1)},\ldots,\omega^{(k)},\ldots,\omega^{(k+m)},V^{(1)},\ldots,V^{(l)},\ldots,V^{(l+n)})\\
&\qquad =T(\omega^{(1)},\ldots,\omega^{(k)},V^{(1)},\ldots,V^{(l)})\\
&\qquad\quad\times S(\omega^{(k+1)},\ldots,\omega^{(k+m)},V^{(l+1)},\ldots,V^{(l+n)})\,.
\tag{1.58}
\end{align*}

注意，$\omega^{(i)}$和$V^{(i)}$是各不相同的对偶矢量和矢量，不是它们的分量。换句话说，先让$T$作用在适当的一组对偶矢量和矢量上，再让$S$作用在其余的上面，然后把两个答案相乘。注意，一般张量积是不可对易的：$T\otimes S\neq S\otimes T$。

现在，取基矢与基对偶矢量的张量积，就可以直截了当地构造出所有$(k,l)$张量构成的空间的基；这个基将由一切形如下式的张量组成

\begin{equation*}
\hat{e}_{(\mu_1)}\otimes\cdots\otimes\hat{e}_{(\mu_k)}
\otimes\hat{\theta}^{(\nu_1)}\otimes\cdots\otimes\hat{\theta}^{(\nu_l)}\,.
\tag{1.59}
\end{equation*}

在四维时空中，这样的基张量总共有$4^{k+l}$个。用分量记号，我们把任意的张量写成

\begin{equation*}
T=T^{\mu_1\cdots\mu_k}{}_{\nu_1\cdots\nu_l}\,
\hat{e}_{(\mu_1)}\otimes\cdots\otimes\hat{e}_{(\mu_k)}
\otimes\hat{\theta}^{(\nu_1)}\otimes\cdots\otimes\hat{\theta}^{(\nu_l)}\,.
\tag{1.60}
\end{equation*}

换个办法，我们也可以通过让张量作用在基矢和基对偶矢量上来定义分量：

\begin{equation*}
T^{\mu_1\cdots\mu_k}{}_{\nu_1\cdots\nu_l}
=T(\hat{\theta}^{(\mu_1)},\ldots,\hat{\theta}^{(\mu_k)},\hat{e}_{(\nu_1)},\ldots,\hat{e}_{(\nu_l)})\,.
\tag{1.61}
\end{equation*}

利用式(1.46)等等，你可以亲自验证这些方程都彼此相容。

与矢量一样，我们通常走捷径，用分量$T^{\mu_1\cdots\mu_k}{}_{\nu_1\cdots\nu_l}$来指代张量$T$。张量作用在一组矢量和一组对偶矢量上的方式，遵循式(1.48)确立的模式：

\begin{equation*}
T(\omega^{(1)},\ldots,\omega^{(k)},V^{(1)},\ldots,V^{(l)})
=T^{\mu_1\cdots\mu_k}{}_{\nu_1\cdots\nu_l}\,
\omega^{(1)}{}_{\mu_1}\cdots\omega^{(k)}{}_{\mu_k}\,
V^{(1)\nu_1}\cdots V^{(l)\nu_l}\,.
\tag{1.62}
\end{equation*}

于是，$(k,l)$张量有$k$个上指标和$l$个下指标。指标的次序显然是重要的，因为张量作用在它的各个宗量上的方式不必相同。

最后，张量分量在洛伦兹变换下的变换，可以借助我们已经知道的基矢和基对偶矢量的变换规律导出。答案正如你按指标位置所预期的：

\begin{equation*}
T^{\mu_1'\cdots\mu_k'}{}_{\nu_1'\cdots\nu_l'}
=\Lambda^{\mu_1'}{}_{\mu_1}\cdots\Lambda^{\mu_k'}{}_{\mu_k}\,
\Lambda^{\nu_1}{}_{\nu_1'}\cdots\Lambda^{\nu_l}{}_{\nu_l'}\,
T^{\mu_1\cdots\mu_k}{}_{\nu_1\cdots\nu_l}\,.
\tag{1.63}
\end{equation*}

这样，每个上指标都像矢量一样变换，每个下指标都像对偶矢量一样变换。

虽然我们把张量定义为从一组矢量和一组对偶矢量到$\mathbf{R}$的线性映射，但没有什么迫使我们非得作用在完整的一组宗量上。于是，$(1,1)$张量也可以作为一个从矢量到矢量的映射来作用：

\begin{equation*}
T^{\mu}{}_{\nu}:\;V^{\nu}\rightarrow T^{\mu}{}_{\nu}V^{\nu}\,.
\tag{1.64}
\end{equation*}

你可以亲自验证$T^{\mu}{}_{\nu}V^{\nu}$是一个矢量（也就是说，服从矢量的变换规律）。类似地，我们可以让一个张量作用在另一个张量上（作用在全部或部分宗量上），得到第三个张量。例如，

\begin{equation*}
U^{\mu}{}_{\nu}=T^{\mu\rho}{}_{\sigma}S^{\sigma}{}_{\rho\nu}
\tag{1.65}
\end{equation*}

就是一个完全合格的$(1,1)$张量。

考虑到这套内容的玄奥性质，你也许会担心我们对张量的介绍有些过于简略。实际上，掌握张量的概念并不需要费很大的力气；无非是把指标摆正而已，而且操作它们的规则是非常自然的。的确，有不少书喜欢把张量\emph{定义}为按式(1.63)变换的一些数的集合。虽然这在操作上是有用的，但它往往掩盖了张量作为几何实体的更深的含义——这些实体具有独立于任何所选坐标系的生命。不过，有一个我们一直含糊带过的微妙之处。对偶矢量、张量、基和线性映射这些概念属于线性代数的领地，只要手头有一个抽象的矢量空间，它们就是适用的。而在我们感兴趣的情形中，我们不只有一个矢量空间，而是时空中每一点都有一个矢量空间。我们更常感兴趣的是张量场，它可以看成时空上的取值为张量的函数。幸运的是，上面定义的那些操作，没有哪一个真正在乎我们处理的是单个矢量空间，还是每个事件一个的矢量空间的汇集。必要时，我们完全可以径直把一些东西称作$x^{\mu}$的函数，而不至于出问题。然而，你应当分清我们引入的这些概念在逻辑上的独立性，以及它们在时空和相对论中的具体应用。

在时空中，我们其实已经见过一些张量的例子，只是没有这样称呼它们。(0,2)张量最熟悉的例子就是度规$\eta_{\mu\nu}$。度规作用在两个矢量上的结果非常有用，所以它有自己的名字，叫做\textbf{内积}（inner product）（或标量积、点积）：

\begin{equation*}
\eta(V,W)=\eta_{\mu\nu}V^{\mu}W^{\nu}=V\cdot W\,.
\tag{1.66}
\end{equation*}

正如在常规的欧几里得点积中一样，我们把内积为零的两个矢量称为\textbf{正交}（orthogonal）。由于内积是标量，它在洛伦兹变换下不变；因此，任何笛卡尔惯性系的基矢——按定义选为正交——在洛伦兹变换之后仍然正交（尽管我们前面注意到它们会“剪拢”到一起）。矢量的\textbf{模}（norm）定义为矢量与自身的内积；与欧几里得空间不同，这个数不是正定的：

\begin{equation*}
\text{若 }\eta_{\mu\nu}V^{\mu}V^{\nu}\text{ 为}\;
\begin{cases}
<0\,, & \ V^{\mu}\text{ 为类时}\,,\\
=0\,, & \ V^{\mu}\text{ 为类光}\,,\\
>0\,, & \ V^{\mu}\text{ 为类空}\,.
\end{cases}
\end{equation*}

（一个矢量可以有零模长而不必是零矢量。）你会注意到，这些术语与我们前面用来分类时空中两点关系的术语相同；当然这绝非偶然，后面我们会更详细地讨论。

另一个张量是型$(1,1)$的\textbf{Kronecker $\delta$ 符号}（Kronecker delta）$\delta^{\mu}{}_{\nu}$。若把它看成从矢量到矢量（或从1-形式到1-形式）的映射，Kronecker $\delta$ 就是恒等映射（identity map）。对于这个独特的张量，我们仿照许多别的参考书的做法，把上下指标放在同一列；讲究纯粹的人也许会写成$\delta^{\mu}{}_{\rho}$或$\delta_{\rho}{}^{\mu}$，但这两者的数值完全相同，所以仅在这个例子上马虎一点不会出问题。

与Kronecker $\delta$和度规相联系的是\textbf{逆度规}（inverse metric）$\eta^{\mu\nu}$，一个型$(2,0)$的张量，（毫不奇怪地）定义为度规的“逆”：

\begin{equation*}
\eta^{\mu\nu}\eta_{\nu\rho}=\eta_{\rho\nu}\eta^{\nu\mu}=\delta^{\mu}{}_{\rho}\,.
\tag{1.67}
\end{equation*}

（它之所以是逆度规，是因为与度规相乘就得到恒等映射。）事实上，你可以验证，逆度规的分量与度规本身的分量完全相同。这只在平直空间的笛卡尔坐标中成立，在更一般的情形下就不再成立。还有\textbf{Levi-Civita符号}（Levi-Civita symbol），一个(0,4)张量：

\begin{equation*}
\tilde{\epsilon}_{\mu\nu\rho\sigma}=
\begin{cases}
+1\,, & \text{若 }\mu\nu\rho\sigma\text{ 是 0123 的偶置换}\\
-1\,, & \text{若 }\mu\nu\rho\sigma\text{ 是 0123 的奇置换}\\
0\,, & \text{其他情形}\,.
\end{cases}
\tag{1.68}
\end{equation*}

这里，“0123 的一个置换”是数0、1、2、3的一个排序，它可以由0123出发交换其中两个数字得到；偶置换通过偶数次这样的交换得到，奇置换通过奇数次得到。例如，$\tilde{\epsilon}_{0321}=-1$。（$\tilde{\epsilon}_{\mu\nu\rho\sigma}$上的波浪号，以及把它称为“符号”而不简单地称为张量，源于这样一个事实：这个对象在更一般的几何或坐标下其实并不是张量；它是一种叫做“张量密度”（tensor density）的东西。定义一个与之相关而是张量的对象是足够直截了当的，我们将把它记作$\epsilon_{\mu\nu\rho\sigma}$，并称之为“Levi-Civita张量”。讨论见第2章。）

上述张量——度规、逆度规、Kronecker $\delta$和Levi-Civita符号——的一个引人注目的性质是：尽管它们都按照张量变换规律(1.63)变换，它们的分量在平直时空的\emph{任何}惯性坐标系中都保持不变。从某种意义上说，这使得它们成为不太典型的张量例子，因为大多数张量并没有这个性质。事实上，尽管我们不去证明，具有这个性质的\emph{只有}这些张量。Kronecker $\delta$还要更加不同寻常，它在任何时空的任何坐标系中都具有完全相同的分量。从张量作为线性映射的定义来看，这很好理解；Kronecker张量可以看成从矢量到矢量（或从对偶矢量到对偶矢量）的恒等映射，显然无论坐标系如何，它都必须具有相同的分量。另一方面，度规及其逆刻画了时空的结构，而Levi-Civita符号则暗地里根本不是真正的张量。因此，当我们放弃平直时空的假设时，就必须更加小心地处理这些对象。

张量的一个更典型的例子是\textbf{电磁场强张量}（electromagnetic field strength tensor）。我们都知道，电磁场由电场矢量$E_i$和磁场矢量$B_i$构成。（记住我们用拉丁指标表示空间分量1、2、3。）实际上，这些只是在空间转动下才是“矢量”，在完整洛伦兹群下则不然。事实上，它们是一个(0,2)张量$F_{\mu\nu}$的分量，其定义为

\begin{equation*}
F_{\mu\nu}=
\begin{pmatrix}
0 & -E_1 & -E_2 & -E_3\\
E_1 & 0 & B_3 & -B_2\\
E_2 & -B_3 & 0 & B_1\\
E_3 & B_2 & -B_1 & 0
\end{pmatrix}
=-F_{\nu\mu}\,.
\tag{1.69}
\end{equation*}

从这个观点出发，通过应用式(1.63)，很容易把一个参考系中的电磁场变换成另一个参考系中的电磁场。张量形式体系的统一威力是显而易见的：与其说是两个矢量的集合，其
关系和变换性质都颇为神秘，我们却用单独一个张量场就描述了全部电磁学。（另一方面，也别得意忘形；有时在单一坐标系中用电场矢量和磁场矢量来工作反而更方便。）

\section{张量的操作}

有了这些例子在手，现在我们可以对张量的某些性质做得更系统一些了。首先考虑\textbf{缩并}（contraction）操作，它把一个 $(k,l)$ 张量变成 $(k-1,l-1)$ 张量。缩并通过对一个上指标和一个下指标求和来进行：
\begin{equation*}
S^{\mu\rho}{}_{\sigma} = T^{\mu\nu\rho}{}_{\sigma\nu}.
\tag{1.70}
\end{equation*}

你可以验证，结果是一个良定义的张量。只有把一个上指标与一个下指标相缩并才是允许的（而不是两个同类型的指标相缩并）；否则结果将\emph{不是}良定义的张量。（所谓良定义的张量，意思要么是“按照张量变换规律进行变换”，要么是“定义了从一组矢量和对偶矢量到实数的唯一的多重线性映射”——任你挑选。）还要注意，指标的顺序是有关系的，因此用不同方式缩并会得到不同的张量；也就是说，在一般情况下
\begin{equation*}
T^{\mu\nu\rho}{}_{\sigma\nu} \neq T^{\mu\rho\nu}{}_{\sigma\nu}
\tag{1.71}
\end{equation*}
成立。

度规和逆度规可以用来对张量\textbf{升指标和降指标}。也就是说，给定一个张量 $T^{\alpha\beta}{}_{\gamma\delta}$，我们可以用度规定义新的张量，并且选择仍用同一个字母 $T$ 来表示它们：
\begin{align*}
T^{\alpha\beta\mu}{}_{\delta} &= \eta^{\mu\gamma}T^{\alpha\beta}{}_{\gamma\delta},\notag\\
T_{\mu}{}^{\beta}{}_{\gamma\delta} &= \eta_{\mu\alpha}T^{\alpha\beta}{}_{\gamma\delta},\notag\\
T_{\mu\nu}{}^{\rho\sigma} &= \eta_{\mu\alpha}\eta_{\nu\beta}\eta^{\rho\gamma}\eta^{\sigma\delta}T^{\alpha\beta}{}_{\gamma\delta},
\tag{1.72}
\end{align*}

如此等等。注意，升降指标并不改变指标相对于其他指标的位置；而且还应注意，自由指标（free indices，不求和的指标）在方程两边必须相同，而哑指标（dummy indices，求和的指标）只出现在一边。举一个例子，我们可以通过升降指标把矢量和对偶矢量互相转换：
\begin{equation*}
\begin{gathered}
V_{\mu} = \eta_{\mu\nu}V^{\nu}\\
\omega^{\mu} = \eta^{\mu\nu}\omega_{\nu}.
\end{gathered}
\tag{1.73}
\end{equation*}

因为度规和逆度规确实互为逆，我们可以自由地把一对正被缩并的指标同时升降：
\begin{equation*}
A^{\lambda}B_{\lambda} = \eta^{\lambda\rho}A_{\rho}\eta_{\lambda\sigma}B^{\sigma} = \delta^{\rho}{}_{\sigma}A_{\rho}B^{\sigma} = A_{\sigma}B^{\sigma}.
\tag{1.74}
\end{equation*}

用度规升降指标的能力解释了，为什么三维平直欧几里得空间中的梯度通常被当成一个普通矢量，尽管我们已经看到它是以对偶矢量的身份出现的；在欧几里得空间中（度规是对角的，且所有元素均为 $+1$），对偶矢量在升指标之后就变成一个分量完全相同的矢量。你可能因此要问：既然如此，为什么我们还要费力去区分它们？一个简单的理由当然是，在洛伦兹时空中两者的分量并不相等：
\begin{equation*}
\omega^{\mu} = (-\omega_{0}, \omega_{1}, \omega_{2}, \omega_{3}).
\tag{1.75}
\end{equation*}

在弯曲时空中，度规的形式一般更为复杂，这种差别要显著得多。但还有一个更深刻的理由，即张量一般都有不依赖于度规的“自然的”定义。尽管我们手头总会有度规可用，但了解我们所引入的每个数学对象的逻辑地位是有益的。梯度凭借其作用于矢量的行为，无须任何度规便有完全良好的定义，而“带上指标的梯度”则不然。（举个例子，我们最终要对泛函作关于度规的变分，因而必须确切知道泛函是如何依赖于度规的，而这一点很容易被指标记号所掩盖。）

继续收集张量行话：如果张量在交换它的某些指标之后保持不变，我们就称它关于这些指标是\textbf{对称的}（symmetric）。于是，如果
\begin{equation*}
S_{\mu\nu\rho} = S_{\nu\mu\rho},
\tag{1.76}
\end{equation*}
我们就说 $S_{\mu\nu\rho}$ 在它的前两个指标上对称；而如果
\begin{equation*}
S_{\mu\nu\rho} = S_{\mu\rho\nu} = S_{\rho\mu\nu} = S_{\nu\mu\rho} = S_{\nu\rho\mu} = S_{\rho\nu\mu},
\tag{1.77}
\end{equation*}
我们就说 $S_{\mu\nu\rho}$ 在它的全部三个指标上都对称。类似地，如果张量在某组指标交换时改变符号，就称它关于这些指标是\textbf{反对称的}（antisymmetric，或称 skew-symmetric）；于是，
\begin{equation*}
A_{\mu\nu\rho} = -A_{\rho\nu\mu}
\tag{1.78}
\end{equation*}
就意味着 $A_{\mu\nu\rho}$ 在它的第一、第三两个指标上反对称（或干脆说“关于 $\mu$ 和 $\rho$ 反对称”）。如果一个张量在它的所有指标上都（反）对称，我们就简称它为（反）对称张量（有时还加上一个画蛇添足的修饰语“完全”）。例如，度规 $\eta_{\mu\nu}$ 和逆度规 $\eta^{\mu\nu}$ 是对称的，而 Levi–Civita 符号 $\tilde{\epsilon}_{\mu\nu\rho\sigma}$ 和电磁场场强张量 $F_{\mu\nu}$ 是反对称的。（请自行验证：如果把一组对称或反对称的指标升降，它们仍保持原样。）注意，上指标与下指标相互交换是没有意义的，所以别经不住诱惑，把 Kronecker $\delta$ 符号 $\delta^{\alpha}{}_{\beta}$ 看成对称的。另一方面，把 $\delta^{\alpha}{}_{\beta}$ 降下一个指标会得到一个对称张量（其实就是度规），这一事实表明指标的顺序实在无关紧要，这正是我们唯独不对这一个张量区分指标位置的原因。

给定任意一个张量，我们可以对它的任意多个上指标或下指标\textbf{对称化}（symmetrize，或反对称化）。对称化就是把相关指标的所有置换求和，再除以项数：
\begin{equation*}
T_{(\mu_1\mu_2\cdots\mu_n)\rho}{}^{\sigma} = \frac{1}{n!}\left(T_{\mu_1\mu_2\cdots\mu_n\rho}{}^{\sigma} + \text{对指标 }\mu_1\cdots\mu_n\text{ 的置换求和}\right),
\tag{1.79}
\end{equation*}
而反对称化则来自交错求和：
\begin{equation*}
T_{[\mu_1\mu_2\cdots\mu_n]\rho}{}^{\sigma} = \frac{1}{n!}\left(T_{\mu_1\mu_2\cdots\mu_n\rho}{}^{\sigma} + \text{对指标 }\mu_1\cdots\mu_n\text{ 的置换交错求和}\right).
\tag{1.80}
\end{equation*}
所谓“交错求和”，是指经奇数次交换得到的置换取负号，即：
\begin{equation*}
T_{[\mu\nu\rho]\sigma} = \frac{1}{6}\left(T_{\mu\nu\rho\sigma} - T_{\mu\rho\nu\sigma} + T_{\rho\mu\nu\sigma} - T_{\nu\mu\rho\sigma} + T_{\nu\rho\mu\sigma} - T_{\rho\nu\mu\sigma}\right).
\tag{1.81}
\end{equation*}
注意，圆括号/方括号表示对称化/反对称化。此外，我们有时想对彼此不相邻的指标作（反）对称化，这时就用竖线标出不参与求和的指标：
\begin{equation*}
T_{(\mu|\nu|\rho)} = \frac{1}{2}\left(T_{\mu\nu\rho} + T_{\rho\nu\mu}\right).
\tag{1.82}
\end{equation*}
如果我们用某个张量上一对对称的上指标去缩并，那么另一个张量的下指标中只有对称的部分会有贡献；于是，
\begin{equation*}
X^{(\mu\nu)}Y_{\mu\nu} = X^{(\mu\nu)}Y_{(\mu\nu)},
\tag{1.83}
\end{equation*}
与 $Y_{\mu\nu}$ 的对称性质无关。（对于反对称的指标，或者一开始对称的就是下指标的情形，也有类似的结论。）对任意\emph{两个}指标，我们可以把张量分解成对称部分和反对称部分，
\begin{equation*}
T_{\mu\nu\rho\sigma} = T_{(\mu\nu)\rho\sigma} + T_{[\mu\nu]\rho\sigma},
\tag{1.84}
\end{equation*}
但这在三个或更多指标的情形一般不再成立，
\begin{equation*}
T_{\mu\nu\rho\sigma} \neq T_{(\mu\nu\rho)\sigma} + T_{[\mu\nu\rho]\sigma},
\tag{1.85}
\end{equation*}
因为存在具有混合对称性的部分，它们既不由对称部分、也不由反对称部分确定。最后，有些人使用的约定是省去 $1/n!$ 因子。这里采用的约定更好，比如说，对称张量就满足
\begin{equation*}
S_{\mu_1\cdots\mu_n} = S_{(\mu_1\cdots\mu_n)},
\tag{1.86}
\end{equation*}
反对称张量亦然。

对一个 $(1,1)$ 型张量 $X^{\mu}{}_{\nu}$，\textbf{迹}（trace）是一个标量，通常省去指标来表示，它其实就是缩并：
\begin{equation*}
X = X^{\lambda}{}_{\lambda}.
\tag{1.87}
\end{equation*}
如果把 $X^{\mu}{}_{\nu}$ 看成一个矩阵，这就是对角分量之和，所以说得通。不过，我们也会在 $(0,2)$ 型张量 $Y_{\mu\nu}$ 的语境下使用迹，这时它的意思是应当先升一个指标（$Y^{\mu}{}_{\nu} = g^{\mu\lambda}Y_{\lambda\nu}$），然后再缩并：
\begin{equation*}
Y = Y^{\lambda}{}_{\lambda} = \eta^{\mu\nu}Y_{\mu\nu}.
\tag{1.88}
\end{equation*}
（必须如此，因为我们不能对两个下指标求和。）虽然这是 $Y^{\mu}{}_{\nu}$ 的对角分量之和，但它显然\emph{不是} $Y_{\mu\nu}$ 的对角分量之和；我们必须先升指标，而这一般会改变分量的数值。例如，你可能以为度规的迹是 $-1+1+1+1=2$，但并不是：
\begin{equation*}
\eta^{\mu\nu}\eta_{\mu\nu} = \delta^{\mu}{}_{\mu} = 4.
\tag{1.89}
\end{equation*}
（在 $n$ 维时 $\delta^{\mu}{}_{\mu}=n$。）没有理由用 $g$（或 $\delta$）来记这个迹，因为它永远是同一个数，即便在我们过渡到度规分量更为复杂的弯曲空间之后也是如此。注意，反对称的 $(0,2)$ 型张量总是无迹的。

迄今为止，我们一直小心地区分两类命题：一类总是成立（在带有任意度规的流形上），另一类只在惯性坐标下的 Minkowski 空间中成立。最重要的区别之一出现在\textbf{偏导数}上。如果在平直时空中采用惯性坐标，那么 $(k,l)$ 张量的偏导数是一个 $(k,l+1)$ 张量；也就是说，
\begin{equation*}
T_{\alpha}{}^{\mu}{}_{\nu} = \partial_{\alpha}R^{\mu}{}_{\nu}
\tag{1.90}
\end{equation*}
在洛伦兹变换下正常变换。然而，在更一般的时空中这不再成立，我们将不得不定义协变导数（covariant derivative）来取代偏导数。尽管如此，在这一特殊情形中，我们仍然可以利用偏导数给出张量这一事实，只要我们保持头脑清醒。[这个警告的唯一例外是标量的偏导数 $\partial_{\alpha}\phi$，它在任何时空中都是一个完全良好的张量（即梯度）。]当然，如果固定某个具体的坐标系，偏导数是一个完全良好的算符，我们会一直使用它；它的失败仅仅在于，它不像我们将要使用的那些张量那样变换（或者说，它所定义的映射不是坐标无关的）。偏导数最有用的性质之一是它们可以对易，
\begin{equation*}
\partial_{\mu}\partial_{\nu}(\cdots) = \partial_{\nu}\partial_{\mu}(\cdots),
\tag{1.91}
\end{equation*}
不论被微分的对象是什么类型。

\section{麦克斯韦方程组}

现在我们已经积累了足够的张量知识，可以用真实的物理来演示其中的一些概念了。具体地说，我们将考察电动力学中的\textbf{麦克斯韦方程组}。用 19 世纪的记号，它们是
\begin{equation*}
\begin{gathered}
\nabla\times\mathbf{B} - \partial_t\mathbf{E} = \mathbf{J}\\
\nabla\cdot\mathbf{E} = \rho\\
\nabla\times\mathbf{E} + \partial_t\mathbf{B} = 0\\
\nabla\cdot\mathbf{B} = 0.
\end{gathered}
\tag{1.92}
\end{equation*}
这里，$\mathbf{E}$ 和 $\mathbf{B}$ 是电场和磁场这两个 3 维矢量，$\mathbf{J}$ 是电流，$\rho$ 是电荷密度，而 $\nabla\times$ 和 $\nabla\cdot$ 是通常的旋度和散度。当然，这些方程在洛伦兹变换下不变；整桩事业正是由此发端的。但它们看上去并不显然不变；我们的张量记号可以弥补这一点。让我们先把它们写成分量记法，
\begin{equation*}
\begin{gathered}
\tilde{\epsilon}^{ijk}\partial_j B_k - \partial_0 E^i = J^i\\
\partial_i E^i = J^0\\
\tilde{\epsilon}^{ijk}\partial_j E_k + \partial_0 B^i = 0\\
\partial_i B^i = 0.
\end{gathered}
\tag{1.93}
\end{equation*}
在这些表达式中，空间指标被随心所欲地升降，全然不去追究度规出现在哪里，因为 $\delta_{ij}$ 就是平直 3 维空间上的度规，$\delta^{ij}$ 是它的逆（作为矩阵两者相等）。因此我们可以随意升降指标，因为分量不会改变。同时，三维 Levi–Civita 符号 $\tilde{\epsilon}^{ijk}$ 的定义与四维的那个完全一样，只是少一个指标（归一化使 $\tilde{\epsilon}^{123}=\tilde{\epsilon}_{123}=1$）。我们用 $J^{0}$ 替换了电荷密度 $\rho$；这是合法的，因为密度和电流合在一起构成了\textbf{电流四矢量}（current 4-vector）$J^{\mu}=(\rho,J^{x},J^{y},J^{z})$。

从式 (1.93) 出发，再加上场强张量 $F_{\mu\nu}$ 的定义式 (1.69)，很容易得到一个完全张量化的、20 世纪版本的麦克斯韦方程组。首先注意到，我们可以用上指标把场强表示为
\begin{equation*}
\begin{gathered}
F^{0i} = E^i\\
F^{ij} = \tilde{\epsilon}^{ijk}B_k.
\end{gathered}
\tag{1.94}
\end{equation*}
为验证这一点，例如可以注意 $F^{01}=\eta^{00}\eta^{11}F_{01}$ 以及 $F^{12}=\tilde{\epsilon}^{123}B_{3}$。于是，式 (1.93) 中的前两个方程变成
\begin{equation*}
\begin{gathered}
\partial_j F^{ij} - \partial_0 F^{0i} = J^i\\
\partial_i F^{0i} = J^0.
\end{gathered}
\tag{1.95}
\end{equation*}
利用 $F^{\mu\nu}$ 的反对称性，我们看到这两式可以合并成单个张量方程
\begin{equation*}
\partial_{\mu}F^{\nu\mu} = J^{\nu}.
\tag{1.96}
\end{equation*}
类似的推理（留作习题）表明，式 (1.93) 中的第三、第四两个方程可以写成
\begin{equation*}
\partial_{[\mu}F_{\nu\lambda]} = 0.
\tag{1.97}
\end{equation*}
很容易验证，$F_{\mu\nu}$ 的反对称性意味着式 (1.97) 可以等价地表示为
\begin{equation*}
\partial_{\mu}F_{\nu\lambda} + \partial_{\nu}F_{\lambda\mu} + \partial_{\lambda}F_{\mu\nu} = 0.
\tag{1.98}
\end{equation*}
于是，四个传统的麦克斯韦方程被两个方程所取代，生动地展示了张量记法的经济性。然而更重要的是，式 (1.96) 和式 (1.97) 的两边都明显地按张量方式变换；因此，只要它们在一个惯性系中成立，它们就必定在任何经洛伦兹变换而来的参考系中成立。这正是张量在相对论中如此有用的原因——我们常常希望表达某种不求助于任何参考系的关系，而方程两边的量在坐标改变时必须以相同的方式变换。按照行话的说法，我们有时会把用张量写出的量称为\textbf{协变的}（covariant）——这与作为“逆变”（contravariant）之反义词的那种“协变”毫无关系。于是我们说，式 (1.96) 和式 (1.97) 合起来构成了麦克斯韦方程组的协变形式，而式 (1.92) 或式 (1.93) 则是非协变的。

\section{能量与动量}

关于照料和饲养张量所要知道的一切，我们如今基本上都过了一遍。下一章我们将更仔细地考察流形和张量的严格定义，但基本的运作机制已经覆盖得相当不错了。在跳到更抽象的数学之前，让我们先回顾一下物理在 Minkowski 时空中是如何运作的。

从单个粒子的世界线讲起。世界线由一个映射 $\mathbf{R}\to M$ 确定，其中 $M$ 是代表时空的流形；我们通常把路径看成参数化曲线 $x^{\mu}(\lambda)$。如前所述，这条路径的切矢量是 $dx^{\mu}/d\lambda$（注意它依赖于参数化方式）。首要感兴趣的对象是切矢量的范数，我们可以用它来刻画路径；如果切矢量在某个参数值 $\lambda$ 处是类时/类光/类空的，我们就说路径在该点是类时/类光/类空的。这解释了为什么同一套字眼既用来给切空间中的矢量分类，也用来给两点之间的间隔分类——因为比如说，连接两个类时分离的点的直线，本身在路径上的每一点都是类时的。

尽管如此，要当心我们在这里耍的一个花招。度规作为一个 $(0,2)$ 型张量，是一台作用于两个矢量（或同一矢量的两份拷贝）而产生一个数的机器。因此，按照范数的符号给切矢量分类是非常自然的。但两点之间的间隔却没有这么自然；它依赖于连接两点的一条特定路径（一条“直线”）的选取，而这种选取又依赖于时空是平直的这一事实（平直性保证了在两点之间可以唯一地选定一条直线）。

让我们从考察一般的路径转到有质量粒子的路径（它们总是类时的）。由于固有时是由沿类时世界线旅行的钟测量的，用 $\tau$ 作为沿路径的参量会很方便。也就是说，我们用式 (1.22) 算出 $\tau(\lambda)$，然后（只要 $\lambda$ 本来就是一个好的参数）把它反解出 $\lambda(\tau)$，此后就可以把路径看成 $x^{\mu}(\tau)$。这种参数化下的切矢量被称为\textbf{四速度}（four-velocity）$U^{\mu}$：
\begin{equation*}
\boxed{U^{\mu} = \frac{dx^{\mu}}{d\tau}.}
\tag{1.99}
\end{equation*}

由于 $d\tau^2 = -\eta_{\mu\nu}dx^{\mu}dx^{\nu}$，四速度自动归一化：
\begin{equation*}
\eta_{\mu\nu}U^{\mu}U^{\nu} = -1.
\tag{1.100}
\end{equation*}

这种绝对的归一化反映出如下事实：四速度并不是穿过空间的速度——后者当然可以取不同的大小——而是“穿过时空的速度”，以它旅行时速率总是一样的。四速度的范数永远是负的，因为我们只对类时轨迹定义它。对类空路径你也可以定义一个类似的矢量；而对类光路径，固有时为零，$\tau$ 不能用作参数，你必须更加小心。在粒子的静止系中，其四速度的分量为 $U^{\mu}=(1,0,0,0)$。

一个与之相关的矢量是\textbf{动量四矢量}（momentum four-vector），定义为
\begin{equation*}
\boxed{p^{\mu} = mU^{\mu},}
\tag{1.101}
\end{equation*}
其中 $m$ 是粒子的质量。质量是一个不依赖于惯性系的固定量，也就是你可能习惯称作“静止质量”的那个东西。事实证明，一劳永逸地把它取作质量，要比认为质量依赖于速度方便得多。粒子的\textbf{能量}就是 $E=p^{0}$，即其动量矢量的类时分量。既然它只是四维矢量的一个分量，它在洛伦兹变换下当然不是不变的；不过这是意料之中的，因为静止粒子的能量与同一粒子运动时的能量本不相同。在粒子的静止系中我们有 $p^{0}=m$；记得我们已令 $c=1$，可见我们已经找到了那个让 Einstein 名声大噪的方程 $E=mc^{2}$。（广义相对论的场方程其实比它更基本，但 $R_{\mu\nu}-\frac{1}{2}Rg_{\mu\nu}=8\pi GT_{\mu\nu}$ 引不起 $E=mc^{2}$ 那种发自肺腑的反应。）在运动的参考系中，我们可以通过施行洛伦兹变换求出 $p^{\mu}$ 的各分量；对于沿 $x$ 轴以三速度 $v=dx/dt$ 运动的粒子，我们有
\begin{equation*}
p^{\mu} = (\gamma m, v\gamma m, 0, 0),
\tag{1.102}
\end{equation*}
其中 $\gamma=1/\sqrt{1-v^{2}}$。当 $v$ 很小时，这给出 $p^{0}=m+\frac{1}{2}mv^{2}$（即我们通常所说的静能加动能）以及 $p^{1}=mv$（即我们通常所说的 Newton 动量）。超出这个近似，我们可以简单地写出
\begin{equation*}
p_{\mu}p^{\mu} = -m^{2},
\tag{1.103}
\end{equation*}
或者
\begin{equation*}
E = \sqrt{m^{2} + \mathbf{p}^{2}},
\tag{1.104}
\end{equation*}
其中 $\mathbf{p}^{2}=\delta_{ij}p^{i}p^{j}$。

相对论以前物理学的核心是 Newton 第二定律，即 $\mathbf{f}=m\mathbf{a}=d\mathbf{p}/dt$。狭义相对论中应当有一个类似的方程，而对方程必须是张量的要求直接引导我们引入一个力四矢量（force four-vector）$f^{\mu}$，它满足
\begin{equation*}
f^{\mu} = m\frac{d^{2}}{d\tau^{2}}x^{\mu}(\tau) = \frac{d}{d\tau}p^{\mu}(\tau).
\tag{1.105}
\end{equation*}
Newton 物理学中最简单的力是引力。但在相对论中，引力并不是用力来描述的，而是用时空自身的曲率。我们转而考虑电磁学。三维的洛伦兹力由 $\mathbf{f}=q(\mathbf{E}+\mathbf{v}\times\mathbf{B})$ 给出，其中 $q$ 是粒子所带的电荷。我们想要这个方程的张量化推广。结果存在一个唯一的答案：
\begin{equation*}
f^{\mu} = -qU^{\lambda}F_{\lambda}{}^{\mu}.
\tag{1.106}
\end{equation*}
你可以自行验证，在小速度极限下它会退化回 Newton 形式。注意，方程必须是张量的这一要求——它同时也是保证洛伦兹不变性的一种方式——严格限制了我们所可能得到的表达式。这是一个非常普遍的现象的例子：表面上看，可能的物理定律五花八门、无穷无尽，而其中一小部分被对称性的要求挑选了出来。

虽然 $p^{\mu}$ 完整描述了单个粒子的能量和动量，我们常常需要处理由数目巨大的粒子组成的延展系统。与其逐一指明每个粒子的动量矢量，不如把系统描述为一种\emph{流体}（fluid）——一种由密度、压强、熵、粘滞性等宏观量刻画的连续体。虽然这样的流体可以由许多具有不同四速度的粒子组成，流体本身却有一个整体的四速度场。想想空气或水这些日常流体就明白了：即便邻近的分子可能具有相当可观的相对速度，对每一个单个的流体元定义一个速度仍然是有意义的。

单独一个动量四矢量场不足以描述流体的能量和动量；我们必须更进一步，定义\textbf{能量-动量张量}（energy-momentum tensor，有时称为能动张量）$T^{\mu\nu}$。这个对称的 $(2,0)$ 型张量告诉我们关于一个系统的能量方面所需要知道的一切：能量密度、压强、应力，等等。$T^{\mu\nu}$ 的一般定义是“四维动量 $p^{\mu}$ 穿过 $x^{\nu}$ 为常数的曲面的通量”。其实，这个定义并不会有多大用处；在第 4 章我们将用作用量对度规的泛函导数来定义能量-动量张量，那将是求 $T^{\mu\nu}$ 显式表达式的更程序化的步骤。不过这里的定义确实提供了某些物理洞见。考虑流体在其静止系中的一个无穷小流体元，此时没有整体流动。那么 $T^{00}$，即“$p^{0}$（能量）在 $x^{0}$（时间）方向上的通量”，正是静止系的\textbf{能量密度}（energy density）$\rho$。类似地，在此系中，$T^{0i}=T^{i0}$ 是动量密度。空间分量 $T^{ij}$ 是动量通量，也就是\emph{应力}（stress）；它们代表流体中相邻无穷小流体元之间的力。$T^{ij}$ 的非对角项代表剪切项，比如由粘滞性引起的那些。像 $T^{11}$ 这样的对角项给出（单位面积上）流体元沿 $x$ 方向所施加的力的 $x$ 分量；这就是我们所说的\textbf{压强}（pressure）的 $x$ 分量 $p_{x}$（不要与动量混淆）。压强有三个分量，在流体静止系（惯性坐标）中由下式给出：
\begin{equation*}
p_i = T^{ii}.
\tag{1.107}
\end{equation*}
这里对 $i$ 不求和。

为了更具体些，让我们从简单的\textbf{尘埃}（dust）例子开始。（宇宙学家往往把“物质”用作尘埃的同义词。）在平直时空中，尘埃可以定义为彼此相对静止的粒子集合。四速度场 $U^{\mu}(x)$ 显然将是各个粒子共同的那个常值四速度。确实，其分量在每一点都相同。定义\textbf{数通量四矢量}（number-flux four-vector）为
\begin{equation*}
N^{\mu} = nU^{\mu},
\tag{1.108}
\end{equation*}
其中 $n$ 是在粒子静止系中测得的粒子数密度。（这听起来不像有坐标不变性，但其实有；在任何参考系中，假如你待在静止系里将会测得的那种数密度是一个固定量。）于是，$N^{0}$ 是在任何其他参考系中测得的粒子数密度，而 $N^{i}$ 是沿 $x^{i}$ 方向的粒子通量。现在设想每个粒子都具有相同的质量 $m$。那么在静止系中，尘埃的能量密度为
\begin{equation*}
\rho = mn.
\tag{1.109}
\end{equation*}
根据定义，能量密度完全确定了尘埃。但 $\rho$ 只量度了静止系中的能量密度；其他参考系呢？我们注意到，$n$ 和 $m$ 都是相应四维矢量在静止系中的 $0$ 分量；具体地，$N^{\mu}=(n,0,0,0)$，$p^{\mu}=(m,0,0,0)$。因此，$\rho$ 正是张量 $p\otimes N$ 在其静止系中 $\mu=0$、$\nu=0$ 的分量。于是我们自然要为尘埃定义能量-动量张量：
\begin{equation*}
T^{\mu\nu}_{\mathrm{dust}} = p^{\mu}N^{\nu} = mn\,U^{\mu}U^{\nu} = \rho\,U^{\mu}U^{\nu},
\tag{1.110}
\end{equation*}
其中 $\rho$ 定义为静止系中的能量密度。（通常你并不是靠这种办法猜出能量-动量张量的，而是从运动方程或作用量原理把它们推导出来。）注意，尘埃沿任何方向的压强都为零；这不足为奇，因为压强来自流体内粒子的无规运动，而我们定义的尘埃恰恰没有这种运动。

尘埃还不够普遍，不足以描述广义相对论中出现的大多数有趣的流体；不过只需稍加推广，就能得到\textbf{理想流体}（perfect fluid）的概念。理想流体是可以完全由两个量确定的流体：静止系能量密度 $\rho$，以及各向同性的静止系压强 $p$。单个参数 $p$ 就规定了每一个方向上的压强。各向同性的一个推论是，$T^{\mu\nu}$ 在其静止系中是对角的——任何一个动量分量在正交方向上都没有净通量。此外，非零的类空分量必定全都相等：$T^{11}=T^{22}=T^{33}$。因此只有两个独立的数：能量密度 $\rho=T^{00}$ 和压强 $p=T^{ii}$；$p$ 不需要下标，因为压强在每一个方向上都相等。于是，理想流体的能量-动量张量在其静止系中取如下形式：
\begin{equation*}
T^{\mu\nu} =
\begin{pmatrix}
\rho & 0 & 0 & 0\\
0 & p & 0 & 0\\
0 & 0 & p & 0\\
0 & 0 & 0 & p
\end{pmatrix}.
\tag{1.111}
\end{equation*}
（记住我们是在平直时空中；引入曲率之后这一点会改变。）当然，我们想要一个在任何参考系中都好用的公式。对尘埃我们有 $T^{\mu\nu}=\rho U^{\mu}U^{\nu}$，于是可以先猜测 $(\rho+p)U^{\mu}U^{\nu}$，它给出
\begin{equation*}
\begin{pmatrix}
\rho+p & 0 & 0 & 0\\
0 & 0 & 0 & 0\\
0 & 0 & 0 & 0\\
0 & 0 & 0 & 0
\end{pmatrix}.
\tag{1.112}
\end{equation*}
老实说，这个猜测并不高明。但从我们想要的结果中减去这个猜测，可以看到还需要加上
\begin{equation*}
\begin{pmatrix}
-p & 0 & 0 & 0\\
0 & p & 0 & 0\\
0 & 0 & p & 0\\
0 & 0 & 0 & p
\end{pmatrix}.
\tag{1.113}
\end{equation*}
幸运的是，它有一个明显的协变推广，即 $p\eta^{\mu\nu}$。于是，理想流体能量-动量张量的一般形式是
\begin{equation*}
\boxed{T^{\mu\nu} = (\rho+p)U^{\mu}U^{\nu} + p\eta^{\mu\nu}.}
\tag{1.114}
\end{equation*}
得出这个公式的过程看上去有些随意，但我们可以对结果抱有充分的信心。既然式 (1.111) 应当是 $T^{\mu\nu}$ 在静止系中的形式，而式 (1.114) 是一个完全张量化的表达式、且在静止系中退化为式 (1.111)，我们就知道式 (1.114) 必定是任何参考系中的正确表达式。

理想流体的概念足够普遍，可以描述五花八门的物质形态。要确定这种流体的演化，我们给定一个把压强和能量密度联系起来的状态方程 $p=p(\rho)$。尘埃是 $p=0$ 的特例，而各向同性的光子气体有 $p=\frac{1}{3}\rho$。一个更奇特的例子是真空能，其能量-动量张量与度规成正比：$T^{\mu\nu}=-\rho_{\mathrm{vac}}\eta^{\mu\nu}$。与式 (1.114) 比较可知，真空能是一种 $p_{\mathrm{vac}}=-\rho_{\mathrm{vac}}$ 的理想流体。在狭义相对论中，“真空中的能量密度”这种概念毫无意义，因为在不含引力的物理里，能量的绝对值并不重要，重要的只是两个态之间的能量差。但在广义相对论中，一切能量都与引力耦合，因此非零真空能的可能性将成为重要的考量，我们将在第 4 章更充分地讨论它。

除了对称之外，$T^{\mu\nu}$ 还有一个更重要的性质：它是\emph{守恒}的。在这里，守恒表现为“散度”为零：
\begin{equation*}
\boxed{\partial_{\mu}T^{\mu\nu} = 0.}
\tag{1.115}
\end{equation*}
这个表达式是一组四个方程，$\nu$ 的每个取值对应一个。$\nu=0$ 的那个方程对应能量守恒，而 $\partial_{\mu}T^{\mu k}=0$ 表示动量第 $k$ 个分量的守恒。让我们把这个方程用于理想流体，此时有
\begin{equation*}
\partial_{\mu}T^{\mu\nu} = \partial_{\mu}(\rho+p)U^{\mu}U^{\nu} + (\rho+p)(U^{\nu}\partial_{\mu}U^{\mu} + U^{\mu}\partial_{\mu}U^{\nu}) + \partial^{\nu}p.
\tag{1.116}
\end{equation*}
为了分析这个方程的含义，一个有益的做法是分别考察：把它投影到沿四速度场 $U^{\mu}$ 的方向和与之正交的方向上，各会得到什么。我们首先注意，归一化条件 $U_{\nu}U^{\nu}=-1$ 蕴含一个有用的恒等式
\begin{equation*}
U_{\nu}\partial_{\mu}U^{\nu} = \frac{1}{2}\partial_{\mu}(U_{\nu}U^{\nu}) = 0.
\tag{1.117}
\end{equation*}
要把式 (1.116) 沿四速度投影，只需把它与 $U_{\nu}$ 缩并：
\begin{equation*}
U_{\nu}\partial_{\mu}T^{\mu\nu} = -\partial_{\mu}(\rho U^{\mu}) - p\,\partial_{\mu}U^{\mu}.
\tag{1.118}
\end{equation*}
令它为零，就得到理想流体相对论性的能量守恒方程。在非相对论极限下它看上去会更眼熟些，在该极限中
\begin{equation*}
U^{\mu} = (1, v^{i}), \qquad |v^{i}| \ll 1, \qquad p \ll \rho.
\tag{1.119}
\end{equation*}
最后一个条件是合理的，因为压强来自单个粒子的无规运动，而在这个极限下这些运动（以及 $U^{\mu}$ 所描写的整体运动）都被视为很小。于是，用通常的非相对论语言，式 (1.118) 变成
\begin{equation*}
\partial_t\rho + \nabla\cdot(\rho\mathbf{v}) = 0,
\tag{1.120}
\end{equation*}
即能量密度的连续性方程。接下来我们考察式 (1.116) 中与四速度正交的部分。要把一个矢量投影到与 $U^{\mu}$ 正交的方向上，我们用投影张量（projection tensor）去乘它：
\begin{equation*}
P^{\sigma}{}_{\nu} = \delta^{\sigma}{}_{\nu} + U^{\sigma}U_{\nu}.
\tag{1.121}
\end{equation*}
为使自己相信它确实管用，可以验证：如果我们有一个平行于 $U^{\mu}$ 的矢量 $V^{\mu}_{\parallel}$，以及另一个垂直于 $U^{\mu}$ 的矢量 $W^{\mu}_{\perp}$，那么投影张量将湮灭平行的矢量而保持正交的矢量：
\begin{equation*}
\begin{gathered}
P^{\sigma}{}_{\nu}V^{\nu}_{\parallel} = 0\\
P^{\sigma}{}_{\nu}W^{\nu}_{\perp} = W^{\sigma}_{\perp}.
\end{gathered}
\tag{1.122}
\end{equation*}
把它用于 $\partial_{\mu}T^{\mu\nu}$，我们得到
\begin{equation*}
P^{\sigma}{}_{\nu}\partial_{\mu}T^{\mu\nu} = (\rho+p)U^{\mu}\partial_{\mu}U^{\sigma} + \partial^{\sigma}p + U^{\sigma}U^{\mu}\partial_{\mu}p.
\tag{1.123}
\end{equation*}
在式 (1.119) 给出的非相对论极限下，令这个表达式的空间分量为零，便得到
\begin{equation*}
\rho\left[\partial_t\mathbf{v} + (\mathbf{v}\cdot\nabla)\mathbf{v}\right] + \nabla p + \mathbf{v}(\partial_t p + \mathbf{v}\cdot\nabla p) = 0.
\tag{1.124}
\end{equation*}

但要注意，最后那一组项中含有 $p$ 的导数与三速度 $\mathbf{v}$ 的乘积，而 $\mathbf{v}$ 已假定为小量；因此与 $\nabla p$ 项相比，这些项可以忽略不计。于是剩下
\begin{equation*}
\rho\left[\partial_t\mathbf{v} + (\mathbf{v}\cdot\nabla)\mathbf{v}\right] = -\nabla p,
\tag{1.125}
\end{equation*}
这正是流体力学中大家熟悉的 Euler 方程。

\section{经典场论}

当我们从狭义相对论过渡到广义相对论时，度规 $\eta_{\mu\nu}$ 将被提升为一个动力学张量场 $g_{\mu\nu}(x)$。广义相对论因此是经典场论的一个具体例子；通过考察定义在平直时空上的经典场，我们可以对这类理论如何运作建立起一些感觉。（我们说经典场论，是相对于量子场论而言的；量子场论是完全不同的另一个故事，我们将在第 9 章对它作简要讨论，但它超出了我们这里的主要兴趣范围。）

让我们从熟悉的一维单粒子经典力学入手，粒子坐标为 $q(t)$。我们可以用“最小作用量原理”导出这种粒子的运动方程：我们寻找某个作用量 $S$ 的临界点（作为轨迹的函数），$S$ 写作
\begin{equation*}
S = \int dt\, L(q, \dot{q}),
\tag{1.126}
\end{equation*}
其中函数 $L(q,\dot{q})$ 就是\textbf{拉格朗日量}（Lagrangian）。点粒子力学中的拉格朗日量通常取如下形式
\begin{equation*}
L = K - V,
\tag{1.127}
\end{equation*}
其中 $K$ 是动能，$V$ 是势能。按照任何一本高等经典力学教材都会讲述的变分法流程，我们可以证明：作用量的临界点［即在微小变分下 $S$ 保持平稳的轨迹 $q(t)$］正是那些满足 \textbf{Euler--Lagrange 方程}（Euler--Lagrange equations）的轨迹，
\begin{equation*}
\frac{\partial L}{\partial q} - \frac{d}{dt}\left(\frac{\partial L}{\partial \dot{q}}\right) = 0.
\tag{1.128}
\end{equation*}
例如，$L = \frac{1}{2}\dot{q}^{2} - V(q)$ 导致
\begin{equation*}
\ddot{q} = -\frac{dV}{dq}.
\tag{1.129}
\end{equation*}

场论也是类似的故事，只不过我们把单个坐标 $q(t)$ 换成了一组依赖于时空的\textbf{场}（fields）$\Phi^{i}(x^{\mu})$，而作用量 $S$ 变成这些场的\emph{泛函}（functional）。泛函不过是以无穷多个变量为自变量的函数，例如某一时空区域中场取的值。泛函常常表示为积分。每个 $\Phi^{i}$ 都是时空上的函数（至少在某个坐标系中是），而 $i$ 是标记我们各个场的指标。例如，在电磁学中（我们将在下面看到），场就是那个叫做“矢势”（vector potential）$A_{\mu}$ 的 1-形式的四个分量：
\begin{equation*}
\Phi^{i} = \{A_{0}, A_{1}, A_{2}, A_{3}\}.
\tag{1.130}
\end{equation*}
我们在这里的做事实在够“俗”的：把一个 1-形式场当成四个不同的函数来想，而不是当成单独的一个张量对象。只要我们固守一个固定的坐标系，这种观点就说得通，而且它会让我们的计算更加直截了当。

在场论中，拉格朗日量可以表示为对一个\textbf{拉格朗日密度}（Lagrange density）$\mathcal{L}$ 的空间积分，$\mathcal{L}$ 是场 $\Phi^{i}$ 及其时空导数 $\partial_{\mu}\Phi^{i}$ 的函数：
\begin{equation*}
L = \int d^{3}x\, \mathcal{L}(\Phi^{i}, \partial_{\mu}\Phi^{i}).
\tag{1.131}
\end{equation*}
于是作用量就是
\begin{equation*}
S = \int dt\, L = \int d^{4}x\, \mathcal{L}(\Phi^{i}, \partial_{\mu}\Phi^{i}).
\tag{1.132}
\end{equation*}
拉格朗日密度是一个洛伦兹标量。我们通常说“拉格朗日量”时，指的其实就是“拉格朗日密度”。定义一个场论时，最方便的做法往往是直接指定拉格朗日密度，所有的运动方程都可以由它轻而易举地导出。

我们将使用“自然单位制”，其中不仅 $c=1$，而且 $\hbar=k=1$，这里 $\hbar=h/2\pi$，$h$ 是 Planck 常数，$k$ 是 Boltzmann 常数。可能有人会反对说，在纯经典的讨论中不该把 $\hbar$ 牵扯进来；但我们在这里做的只是选取单位，并不是在决定物理。（如果我们要把场论量子化并得到粒子，$\hbar$ 的相关性就会显现出来，不过我们现在不会走到那一步。）在自然单位制中我们有
\begin{equation*}
[\text{能量}] = [\text{质量}] = [(\text{长度})^{-1}] = [(\text{时间})^{-1}].
\tag{1.133}
\end{equation*}
我们将最经常用能量或质量作为基本单位。既然作用量是 $L$（量纲为能量）对时间的积分，它就是无量纲的：
\begin{equation*}
[S] = [E][T] = M^{0}.
\tag{1.134}
\end{equation*}
体积元的量纲为
\begin{equation*}
[d^{4}x] = M^{-4},
\tag{1.135}
\end{equation*}
因此，为了使作用量无量纲，我们要求拉格朗日密度具有量纲
\begin{equation*}
[\mathcal{L}] = M^{4}.
\tag{1.136}
\end{equation*}

Euler--Lagrange 方程来自这样的要求：作用量在场的微小变分下不变，
\begin{equation*}
\Phi^{i} \rightarrow \Phi^{i} + \delta\Phi^{i},
\tag{1.137}
\end{equation*}
\begin{equation*}
\partial_{\mu}\Phi^{i} \rightarrow \partial_{\mu}\Phi^{i} + \delta(\partial_{\mu}\Phi^{i}) = \partial_{\mu}\Phi^{i} + \partial_{\mu}(\delta\Phi^{i}).
\tag{1.138}
\end{equation*}
$\partial_{\mu}\Phi^{i}$ 的变分表达式不过是 $\Phi^{i}$ 的变分的导数。既然假定 $\delta\Phi^{i}$ 是小量，我们就可以在这个变分下对拉格朗日密度作泰勒展开：
\begin{align*}
\mathcal{L}(\Phi^{i}, \partial_{\mu}\Phi^{i}) &\rightarrow \mathcal{L}(\Phi^{i} + \delta\Phi^{i}, \partial_{\mu}\Phi^{i} + \partial_{\mu}\delta\Phi^{i})\notag\\
&= \mathcal{L}(\Phi^{i}, \partial_{\mu}\Phi^{i}) + \frac{\partial\mathcal{L}}{\partial\Phi^{i}}\delta\Phi^{i} + \frac{\partial\mathcal{L}}{\partial(\partial_{\mu}\Phi^{i})}\partial_{\mu}(\delta\Phi^{i}).
\tag{1.139}
\end{align*}
相应地，作用量变为 $S\rightarrow S+\delta S$，其中
\begin{equation*}
\delta S = \int d^{4}x\left[\frac{\partial\mathcal{L}}{\partial\Phi^{i}}\delta\Phi^{i} + \frac{\partial\mathcal{L}}{\partial(\partial_{\mu}\Phi^{i})}\partial_{\mu}(\delta\Phi^{i})\right].
\tag{1.140}
\end{equation*}
我们希望通过对第二项作分部积分，把 $\delta\Phi^{i}$ 从被积函数中提出来：
\begin{align*}
\int d^{4}x\, \frac{\partial\mathcal{L}}{\partial(\partial_{\mu}\Phi^{i})}\partial_{\mu}(\delta\Phi^{i}) &= -\int d^{4}x\, \partial_{\mu}\left(\frac{\partial\mathcal{L}}{\partial(\partial_{\mu}\Phi^{i})}\right)\delta\Phi^{i}\notag\\
&\quad + \int d^{4}x\, \partial_{\mu}\left(\frac{\partial\mathcal{L}}{\partial(\partial_{\mu}\Phi^{i})}\delta\Phi^{i}\right).
\tag{1.141}
\end{align*}
最后一项是一个全导数——形如 $\partial_{\mu}V^{\mu}$ 的某个量的积分——借助 Stokes 定理（是四维版本；讨论见附录 E）可以把它化为一个表面项。既然我们处理的是变分问题，我们可以选择考虑那些在边界上（连同其导数）为零的变分。因此在这种情形下，传统上总是肆无忌惮地作分部积分，把边界贡献一概略去。（有时这样做是不行的，例如 Yang--Mills 理论中的瞬子计算。）

于是我们剩下
\begin{equation*}
\delta S = \int d^{4}x\left[\frac{\partial\mathcal{L}}{\partial\Phi^{i}} - \partial_{\mu}\left(\frac{\partial\mathcal{L}}{\partial(\partial_{\mu}\Phi^{i})}\right)\right]\delta\Phi^{i}.
\tag{1.142}
\end{equation*}
泛函 $S$ 相对于函数 $\Phi^{i}$ 的\textbf{泛函导数}（functional derivative）$\delta S/\delta\Phi^{i}$ 定义为满足
\begin{equation*}
\delta S = \int d^{4}x\, \frac{\delta S}{\delta\Phi^{i}}\delta\Phi^{i},
\tag{1.143}
\end{equation*}
只要这样的表达式成立。因此我们可以把“$S$ 处于临界点”这一说法表述为泛函导数为零。于是，我们这个场论的最终运动方程就是：
\begin{equation*}
\boxed{\;\frac{\delta S}{\delta\Phi^{i}} = \frac{\partial\mathcal{L}}{\partial\Phi^{i}} - \partial_{\mu}\left(\frac{\partial\mathcal{L}}{\partial(\partial_{\mu}\Phi^{i})}\right) = 0.\;}
\tag{1.144}
\end{equation*}
这些就是所谓的平直时空中场论的 Euler--Lagrange 方程。

场最简单的例子是实标量场：
\begin{equation*}
\phi(x^{\mu}) : (\text{时空}) \rightarrow \mathbf{R}.
\tag{1.145}
\end{equation*}
稍微复杂一些的例子可以包括复标量场，或者从时空到任意矢量空间、乃至任意流形的映射（有时称为“非线性 $\sigma$ 模型”（nonlinear sigma models））。量子化之后，场的激发作为粒子是可以观测到的。标量场给出无自旋粒子，而矢量场和其他张量给出更高自旋的粒子。如果场是复的而非实的，它将有两个而非仅仅一个自由度，这会被诠释为一个粒子加上一个不同的反粒子。实场就是它们自己的反粒子。实标量场的一个例子是中性的 $\pi$ 介子。

那么我们来考虑单个实标量场的经典力学。它将具有一个能量密度，是时空的局域函数，并包含各种贡献：
\begin{equation*}
\begin{array}{lll}
\text{动能} &:& \frac{1}{2}\dot{\phi}^{2}\\
\text{梯度能量} &:& \frac{1}{2}(\nabla\phi)^{2}\\
\text{势能} &:& V(\phi).
\end{array}
\tag{1.146}
\end{equation*}
实际上，虽然势是洛伦兹不变的函数，但动能和梯度能量本身并不是洛伦兹不变的；不过我们可以把它们组合成一个明显洛伦兹不变的形式：
\begin{equation*}
-\frac{1}{2}\eta^{\mu\nu}(\partial_{\mu}\phi)(\partial_{\nu}\phi) = \frac{1}{2}\dot{\phi}^{2} - \frac{1}{2}(\nabla\phi)^{2}.
\tag{1.147}
\end{equation*}
［组合 $\eta^{\mu\nu}(\partial_{\mu}\phi)(\partial_{\nu}\phi)$ 常被缩写为 $(\partial\phi)^{2}$。］于是，对我们的单个实标量场，类比于点粒子情形的 $L=K-V$，一个合理的拉格朗日量选择是
\begin{equation*}
\mathcal{L} = -\frac{1}{2}\eta^{\mu\nu}(\partial_{\mu}\phi)(\partial_{\nu}\phi) - V(\phi).
\tag{1.148}
\end{equation*}
这是把“动能减势能”推广成“动能减梯度能量减势能密度”。注意，由于 $[\mathcal{L}]=M^{4}$，必须有 $[V]=M^{4}$。又由于 $[\partial_{\mu}]=[\partial/\partial x^{\mu}]=M^{1}$，我们有
\begin{equation*}
[\phi] = M^{1}.
\tag{1.149}
\end{equation*}

对于拉格朗日量 (1.148)，我们有
\begin{equation*}
\frac{\partial\mathcal{L}}{\partial\phi} = -\frac{dV}{d\phi}, \qquad \frac{\partial\mathcal{L}}{\partial(\partial_{\mu}\phi)} = -\eta^{\mu\nu}\partial_{\nu}\phi.
\tag{1.150}
\end{equation*}
这两个式子中的第二个有点棘手，让我们慢慢来做。在对拉格朗日量求导时，诀窍是确保指标的摆放“相容”（也就是说，如果求导所针对的对象上带一个下指标，那么当同类对象出现在被求导的东西里时，它也应当只带下指标），同时还要保证指标严格不同。第一条在我们的例子中已经满足，因为我们是在把 $\partial_{\mu}\phi$ 的函数对 $\partial_{\mu}\phi$ 求导。往后我们就需要更加小心了。为了满足第二条，我们只需重新命名哑指标：
\begin{equation*}
\eta^{\mu\nu}(\partial_{\mu}\phi)(\partial_{\nu}\phi) = \eta^{\rho\sigma}(\partial_{\rho}\phi)(\partial_{\sigma}\phi).
\tag{1.151}
\end{equation*}
然后我们可以使用如下的一般规则：对任何带一个指标的对象，比如 $V_{\mu}$，有
\begin{equation*}
\frac{\partial V_{\alpha}}{\partial V_{\beta}} = \delta^{\beta}{}_{\alpha}
\tag{1.152}
\end{equation*}
因为 $V_{\alpha}$ 的每个分量都被当作一个独立的变量。于是我们有
\begin{align*}
\frac{\partial}{\partial(\partial_{\mu}\phi)}\left[\eta^{\rho\sigma}(\partial_{\rho}\phi)(\partial_{\sigma}\phi)\right] &= \eta^{\rho\sigma}\left[\delta^{\mu}{}_{\rho}(\partial_{\sigma}\phi) + (\partial_{\rho}\phi)\delta^{\mu}{}_{\sigma}\right]\notag\\
&= \eta^{\mu\sigma}(\partial_{\sigma}\phi) + \eta^{\rho\mu}(\partial_{\rho}\phi) = 2\eta^{\mu\nu}\partial_{\nu}\phi.
\tag{1.153}
\end{align*}
这就得出了 (1.150) 中的第二个表达式。

把 (1.150) 代入 (1.144)，就得到运动方程
\begin{equation*}
\Box\phi - \frac{dV}{d\phi} = 0,
\tag{1.154}
\end{equation*}
其中 $\Box\equiv\eta^{\mu\nu}\partial_{\mu}\partial_{\nu}$ 被称为 \textbf{d'Alembert 算符}（d'Alembertian）。注意，这个方程体现了我们的度规号差约定 $(-+++)$；如果采用另一种 $(+---)$ 约定，符号就会反过来。在平直时空中，(1.154) 等价于
\begin{equation*}
\ddot{\phi} - \nabla^{2}\phi + \frac{dV}{d\phi} = 0.
\tag{1.155}
\end{equation*}

势 $V$ 的一个流行选择是简单谐振子的势，$V(\phi)=\frac{1}{2}m^{2}\phi^{2}$。参数 $m$ 叫做场的质量，你应该会注意到量纲恰好正确。你可能觉得奇怪，场怎么会有质量。当我们把场量子化，会发现动量本征态是粒子的集合，每个粒子的质量为 $m$。在经典层次上，我们把“质量”看作对场动力学的一种方便的刻画。于是我们的运动方程就是
\begin{equation*}
\Box\phi - m^{2}\phi = 0,
\tag{1.156}
\end{equation*}
这就是著名的 \textbf{Klein--Gordon 方程}。这是一个线性微分方程，所以两个解的和仍是解；一组完备的解（平面波形式）很容易找到，你可以自行核验。

电磁学提供了场论的一个稍微复杂些的例子。我们提到过，相关的场是\textbf{矢势}（vector potential）$A_{\mu}$；其类时分量 $A_{0}$ 可以等同于静电势 $\Phi$，类空分量则等同于传统的矢量势 $\mathbf{A}$（磁场由 $\mathbf{B}=\nabla\times\mathbf{A}$ 给出）。场强张量（其分量由式 (1.69) 给出）与矢势的关系为
\begin{equation*}
F_{\mu\nu} = \partial_{\mu}A_{\nu} - \partial_{\nu}A_{\mu}.
\tag{1.157}
\end{equation*}
从这个定义我们看到，场强张量具有重要的\textbf{规范不变性}（gauge invariance）性质：当我们对矢势施加一个\textbf{规范变换}（gauge transformation）
\begin{equation*}
A_{\mu} \rightarrow A_{\mu} + \partial_{\mu}\lambda(x),
\tag{1.158}
\end{equation*}
时，场强张量保持不变：
\begin{equation*}
F_{\mu\nu} \rightarrow F_{\mu\nu} + \partial_{\mu}\partial_{\nu}\lambda - \partial_{\nu}\partial_{\mu}\lambda = F_{\mu\nu}.
\tag{1.159}
\end{equation*}
最后一个等号源于偏导数可以对易这一事实：$\partial_{\mu}\partial_{\nu}=\partial_{\nu}\partial_{\mu}$。规范不变性是我们理解电磁学的一种根本性对称，所有可观测量都必须是规范不变的。因此，虽然这个理论的动力学场（我们对它变分作用量来导出运动方程）是 $A_{\mu}$，物理量一般却会用 $F_{\mu\nu}$ 来表达。

我们已经知道，电磁学的动力学方程就是麦克斯韦方程组，即式 (1.96) 和式 (1.97)。给定场强张量用矢势表达的定义，(1.97) 实际上是自动满足的：
\begin{equation*}
\partial_{[\mu}F_{\nu\sigma]} = \partial_{[\mu}\partial_{\nu}A_{\sigma]} - \partial_{[\mu}\partial_{\sigma}A_{\nu]} = 0,
\tag{1.160}
\end{equation*}
同样是因为偏导数可以对易。另一方面，(1.96) 则等价于如下形式的 Euler--Lagrange 方程
\begin{equation*}
\frac{\partial\mathcal{L}}{\partial A_{\nu}} - \partial_{\mu}\left(\frac{\partial\mathcal{L}}{\partial(\partial_{\mu}A_{\nu})}\right) = 0,
\tag{1.161}
\end{equation*}
只要我们有先见之明地选取拉格朗日量为
\begin{equation*}
\mathcal{L} = -\frac{1}{4}F_{\mu\nu}F^{\mu\nu} + A_{\mu}J^{\mu}.
\tag{1.162}
\end{equation*}

对这个选择，Euler--Lagrange 方程中的第一项是直截了当的：
\begin{equation*}
\frac{\partial\mathcal{L}}{\partial A_{\nu}} = \delta^{\nu}{}_{\mu}J^{\mu} = J^{\nu}.
\tag{1.163}
\end{equation*}
第二项更棘手些。首先我们把 $F_{\mu\nu}F^{\mu\nu}$ 写成
\begin{equation*}
F_{\mu\nu}F^{\mu\nu} = F_{\alpha\beta}F^{\alpha\beta} = \eta^{\alpha\rho}\eta^{\beta\sigma}F_{\alpha\beta}F_{\rho\sigma}.
\tag{1.164}
\end{equation*}
我们想在 $F_{\mu\nu}$ 上使用下指标，因为我们是相对于具有下指标的 $\partial_{\mu}A_{\nu}$ 求导。同样，我们把 $F_{\mu\nu}F^{\mu\nu}$ 的哑指标也换掉，因为我们希望被求导的对象与求导所针对的对象带有不同的指标。一旦你熟悉了这一套，它就会成为你的第二天性，你将不再需要这么多步骤。这使我们能够写下
\begin{equation*}
\frac{\partial(F_{\alpha\beta}F^{\alpha\beta})}{\partial(\partial_{\mu}A_{\nu})} = \eta^{\alpha\rho}\eta^{\beta\sigma}\left[\left(\frac{\partial F_{\alpha\beta}}{\partial(\partial_{\mu}A_{\nu})}\right)F_{\rho\sigma} + F_{\alpha\beta}\left(\frac{\partial F_{\rho\sigma}}{\partial(\partial_{\mu}A_{\nu})}\right)\right].
\tag{1.165}
\end{equation*}
然后，由于 $F_{\alpha\beta}=\partial_{\alpha}A_{\beta}-\partial_{\beta}A_{\alpha}$，我们有
\begin{equation*}
\frac{\partial F_{\alpha\beta}}{\partial(\partial_{\mu}A_{\nu})} = \delta^{\mu}{}_{\alpha}\delta^{\nu}{}_{\beta} - \delta^{\mu}{}_{\beta}\delta^{\nu}{}_{\alpha}.
\tag{1.166}
\end{equation*}
把 (1.166) 与 (1.165) 结合起来就得到
\begin{align*}
\frac{\partial(F_{\alpha\beta}F^{\alpha\beta})}{\partial(\partial_{\mu}A_{\nu})} &= \eta^{\alpha\rho}\eta^{\beta\sigma}\left[\left(\delta^{\mu}{}_{\alpha}\delta^{\nu}{}_{\beta} - \delta^{\mu}{}_{\beta}\delta^{\nu}{}_{\alpha}\right)F_{\rho\sigma} + \left(\delta^{\mu}{}_{\rho}\delta^{\nu}{}_{\sigma} - \delta^{\mu}{}_{\sigma}\delta^{\nu}{}_{\rho}\right)F_{\alpha\beta}\right]\notag\\
&= \left(\eta^{\mu\rho}\eta^{\nu\sigma} - \eta^{\nu\rho}\eta^{\mu\sigma}\right)F_{\rho\sigma} + \left(\eta^{\alpha\mu}\eta^{\beta\nu} - \eta^{\alpha\nu}\eta^{\beta\mu}\right)F_{\alpha\beta}\notag\\
&= F^{\mu\nu} - F^{\nu\mu} + F^{\mu\nu} - F^{\nu\mu}\notag\\
&= 4F^{\mu\nu},
\tag{1.167}
\end{align*}
因此
\begin{equation*}
\frac{\partial\mathcal{L}}{\partial(\partial_{\mu}A_{\nu})} = -F^{\mu\nu}.
\tag{1.168}
\end{equation*}
再把 (1.163) 和 (1.168) 代入 (1.161)，恰好给出 (1.96)：
\begin{equation*}
\partial_{\mu}F^{\nu\mu} = J^{\nu}.
\tag{1.169}
\end{equation*}
注意，我们交换了 $F^{\mu\nu}$ 上指标的顺序，免得自己吃进一个讨厌的负号。

你可能纳闷，如果我们能够在知道拉格朗日量之前就发明运动方程（正如 Maxwell 对他的方程所做的那样），那么引入拉格朗日表述的目的何在。理由有好几个：首先是基本的简单性——设定一个单一的时空标量函数，即拉格朗日密度，而不是一堆（可能是张量值的）运动方程。另一个理由是实施对称性十分便利；要求作用量在某一对称性下不变，就保证了动力学同样尊重这一对称性。最后，正如我们将在第 4 章看到的，作用量经由一个直接的手续（包括对度规本身求变分）导向一个唯一的能量-动量张量。把这个手续应用于 (1.148)，就直接得到标量场论的能量-动量张量，
\begin{equation*}
T^{\mu\nu}_{\text{scalar}} = \eta^{\mu\lambda}\eta^{\nu\sigma}\partial_{\lambda}\phi\,\partial_{\sigma}\phi - \eta^{\mu\nu}\left[\frac{1}{2}\eta^{\lambda\sigma}\partial_{\lambda}\phi\,\partial_{\sigma}\phi + V(\phi)\right].
\tag{1.170}
\end{equation*}
类似地，从 (1.162) 出发我们可以导出电磁学的能量-动量张量，
\begin{equation*}
T^{\mu\nu}_{\text{EM}} = F^{\mu\lambda}F^{\nu}{}_{\lambda} - \frac{1}{4}\eta^{\mu\nu}F^{\lambda\sigma}F_{\lambda\sigma}.
\tag{1.171}
\end{equation*}
利用相应的运动方程，你可以证明这些能量-动量张量是守恒的，$\partial_{\mu}T^{\mu\nu}=0$（习题里也会要求你这样做）。

我们考虑过的这两个例子——标量场论和电磁学——是我们当前对自然的理解中大部分内容的范本。粒子物理学的标准模型由三类场组成：规范场、Higgs 场和费米子。规范场描述自然的“力”，除电磁相互作用之外还包括强核力和弱核力。产生核力的那些规范场与电磁学中一样，由 1-形式势来描述；区别在于它们是矩阵值的，而非普通的 1-形式，因此与规范变换相对应的对称群是非对易（非阿贝尔）的对称性。Higgs 场和我们描述过的标量场非常相像，尽管它们也是矩阵值的。费米子包括轻子（例如电子和中微子）和夸克，它们不由我们这里讨论过的任何张量场来描述，而是由一种不同类型的场——所谓\textbf{旋量}（spinor）——来描述。在本书中我们不会展开讨论旋量，但它们在粒子物理中扮演着关键的角色，而且它们与引力的耦合既有趣又微妙。量子化之后，这些场给出具有不同自旋的粒子：规范场是自旋-1 的，标量场是自旋-0 的，而标准模型的费米子是自旋-$\frac{1}{2}$ 的。

在结束本章之前，让我们问一个简单得让人不好意思的问题：为什么我们要考虑这一个经典场论，而不是别的某个？说得更具体些，假设我们在自然界中发现了某种粒子，并且我们知道想用哪一类场来描述它；那么我们该如何为这个场挑选拉格朗日量？例如，当我们写下标量场拉格朗日量 (1.148) 时，为什么我们没有包含形如下式的项
\begin{equation*}
\mathcal{L}' = \lambda\phi^{2}\eta^{\mu\nu}(\partial_{\mu}\phi)(\partial_{\nu}\phi),
\tag{1.172}
\end{equation*}
其中 $\lambda$ 是一个耦合常数？归根结底，我们当然是通过试错来工作，努力拟合实验提供给我们的数据。在经典场论中，我们能做的也就不多了；一般来说我们会从一个简单的拉格朗日量出发，如果第一次尝试与数据不符，或许再把它复杂化。但量子场论实际上提供了一些简单的指导原则，而且既然我们把经典场论当作某个底层量子理论的近似来使用，利用这些原则就是有意义的。长话短说：量子场论允许任意高能量的“虚”过程对我们低能下观测到的东西作出贡献。所幸的是，这些过程的效应可以总结在一个低能\textbf{有效场论}（effective field theory）之中。在有效理论——这才是我们实际观测到的理论——中，高能过程的结果只不过是把我们理论的耦合常数“重整化”。考虑一个任意的耦合常数，它可以表示成参数 $\mu$（量纲为质量）的某次幂，$\lambda=\mu^{d}$（除非 $\lambda$ 无量纲，那样的话讨论会变得更微妙）。非常粗略地说，\emph{高能过程的效应将是使 $\mu$ 变得非常大。}稍微具体一点：$\mu$ 将被推高到一个新物理开始介入的尺度，不管那会是什么。因此，我们可能想往拉格朗日量里添加的潜在高阶项都被压低了，因为它们所乘的耦合常数非常小。例如对 (1.172) 来说，必须有 $\lambda=\mu^{-2}$，于是 $\lambda$ 会微乎其微（因为 $\mu$ 会很大）。我们在拉格朗日量中能放进去的最低阶项才带有无量纲的耦合常数（或者量纲为质量的正幂次的耦合常数），所以在低能下我们只需要操心这些项。场论的这一特点带来了巨大的简化，使我们得以考察所有想要研究的模型。

正如本节开头所提到的，广义相对论本身就是一个经典场论，其动力学场就是度规张量。然而，把广义相对论看作某种与众不同的东西也是公允的；其他经典场论在很大程度上依赖于一个预先存在的时空几何，而在广义相对论中，几何是由运动方程决定的。（这一想法也有例外，称为拓扑场论，在其中度规根本不露面。）接下来几章里我们的任务是探究由时空度规所刻画的弯曲几何的性质，然后在第 4 章把这些概念付诸实践，构造出一个引力理论。

\section{习题}

% 习题编号照原书
\textbf{1.}\quad 考虑一个惯性系 $S$，其坐标为 $x^{\mu}=(t,x,y,z)$，以及一个参考系 $S'$，其坐标 $x^{\mu'}$ 与 $S$ 之间通过一个沿 $y$ 轴方向、速度参量为 $v$ 的推动相联系。设想有一堵墙静止在 $S'$ 中，位于直线 $x'=-y'$ 上。从 $S$ 的观点来看，一个击中墙壁的球（在 $x$--$y$ 平面内运动）的入射角与反射角之间是什么关系？碰前和碰后的速度又如何？

\textbf{2.}\quad 设想空间（不是时空）实际上是一个有限的盒子，或者用更考究的说法，一个大小为 $L$ 的三维环面（three-torus）。我们的意思是：存在一个坐标系 $x^{\mu}=(t,x,y,z)$，使得坐标为 $(t,x,y,z)$ 的每个点都与坐标为 $(t,x+L,y,z)$、$(t,x,y+L,z)$ 和 $(t,x,y,z+L)$ 的每个点相\emph{认同}（identified）。注意时间坐标是相同的。现在考虑两个观者：观者 $A$ 在这个坐标系中静止（空间坐标为常数），而观者 $B$ 以恒定速度 $v$ 沿 $x$ 方向运动。$A$ 和 $B$ 从同一事件出发，当 $A$ 保持不动时，$B$ 绕宇宙一圈回来，与 $A$ 的世界线相交而完全不必加速（因为宇宙是周期性的）。在这段间隔中 $A$ 和 $B$ 经历的相对固有时是多少？这与你对洛伦兹不变性的理解一致吗？

\textbf{3.}\quad 三个事件 $A$、$B$、$C$，在观者 $\mathcal{O}$ 看来按 $ABC$ 的顺序发生。另一位观者 $\bar{\mathcal{O}}$ 看到这些事件按 $CBA$ 的顺序发生。是否可能有第三位观者看到这些事件按 $ACB$ 的顺序发生？请通过画一张时空图来支持你的结论。

\textbf{4.}\quad 投影效应会骗你以为某个天体物理对象在“超光速地”运动。考虑一个类星体，它以速率 $v$、与观者视线成 $\theta$ 角抛出气体。投影到天空上，气体看起来垂直于视线方向行进，角速度为 $v_{\mathrm{app}}/D$，其中 $D$ 是到类星体的距离，$v_{\mathrm{app}}$ 是表观速度。推导用 $v$ 和 $\theta$ 表出 $v_{\mathrm{app}}$ 的表达式。证明，对适当的 $v$ 和 $\theta$ 值，$v_{\mathrm{app}}$ 可以大于 1。

\textbf{5.}\quad 粒子物理学家太习惯于取 $c=1$ 了，以至于他们用能量的单位来量度质量。特别是，他们倾向于使用电子伏特（$1\,\text{eV} = 1.6\times 10^{-12}\,\text{erg} = 1.8\times 10^{-33}\,\text{g}$），或者更常用的 keV、MeV 和 GeV（分别为 $10^{3}\,\text{eV}$、$10^{6}\,\text{eV}$ 和 $10^{9}\,\text{eV}$）。$\mu$ 子已被测得质量为 $0.106\,\text{GeV}$，静止系寿命为 $2.19\times 10^{-6}$ 秒。设想这样一个 $\mu$ 子在粒子加速器直径 1 千米的圆形储存环中运动，其总能量为 $1000\,\text{GeV}$。从实验者的观点来看，它看起来能活多久？它会绕环走过多少弧度？

\textbf{6.}\quad 在欧几里得三维空间中，设 $p$ 是坐标为 $(x,y,z)=(1,0,-1)$ 的点。考虑下列经过 $p$ 的曲线：
\begin{gather*}
x^{i}(\lambda) = \left(\lambda, (\lambda-1)^{2}, -\lambda\right)\\
x^{i}(\mu) = \left(\cos\mu, \sin\mu, \mu-1\right)\\
x^{i}(\sigma) = \left(\sigma^{2}, \sigma^{3}+\sigma^{2}, \sigma\right).
\end{gather*}
\textbf{(a)}\quad 计算这些曲线在 $p$ 点的切矢量在坐标基 $\{\partial_{x}, \partial_{y}, \partial_{z}\}$ 下的分量。

\textbf{(b)}\quad 设 $f = x^{2}+y^{2}-yz$。计算 $df/d\lambda$、$df/d\mu$ 和 $df/d\sigma$。

\textbf{7.}\quad 设想我们有一个张量 $X^{\mu\nu}$ 和一个矢量 $V^{\mu}$，分量为
\begin{equation*}
X^{\mu\nu} =
\begin{pmatrix}
2 & 0 & 1 & -1\\
-1 & 0 & 3 & 2\\
-1 & 1 & 0 & 0\\
-2 & 1 & 1 & -2
\end{pmatrix},
\qquad
V^{\mu} = (-1, 2, 0, -2).
\end{equation*}
求下列各量的分量：

\textbf{(a)}\quad $X^{\mu}{}_{\nu}$

\textbf{(b)}\quad $X_{\mu}{}^{\nu}$

\textbf{(c)}\quad $X^{(\mu\nu)}$

\textbf{(d)}\quad $X_{[\mu\nu]}$

\textbf{(e)}\quad $X^{\lambda}{}_{\lambda}$

\textbf{(f)}\quad $V^{\mu}V_{\mu}$

\textbf{(g)}\quad $V_{\mu}X^{\mu\nu}$

\textbf{8.}\quad 若 $\partial_{\nu}T^{\mu\nu}=Q^{\mu}$，那么空间矢量 $Q^{i}$ 在物理上代表什么？请用尘埃的能量-动量张量来论证你的回答。

\textbf{9.}\quad 对于一个分立点粒子系统，能量-动量张量取如下形式
\begin{equation*}
T_{\mu\nu} = \sum_{a} \frac{p^{(a)}_{\mu}\,p^{(a)}_{\nu}}{p^{0(a)}}\,\delta^{(3)}(\mathbf{x}-\mathbf{x}^{(a)}),
\tag{1.173}
\end{equation*}
其中指标 $a$ 标记不同的粒子。证明，对于速度作各向同性分布的稠密粒子集合，我们可以把单个粒子的世界线抹平，从而得到理想流体的能量-动量张量 (1.114)。

\textbf{10.}\quad 利用作用到 $F_{\mu\nu}$ 上的张量变换规律，证明电场和磁场的三维矢量 $\mathbf{E}$ 和 $\mathbf{B}$ 在下列情形下如何变换：

\textbf{(a)}\quad 绕 $y$ 轴的转动；

\textbf{(b)}\quad 沿 $z$ 轴的推动。

\textbf{11.}\quad 验证 (1.98) 的确等价于 (1.97)，并且两者都等价于 (1.93) 中的最后两个方程。

\textbf{12.}\quad 考虑我们明确讨论过的两个场论：Maxwell 电磁学（令 $J^{\mu}=0$）以及由 (1.148) 定义的标量场论。

\textbf{(a)}\quad 用三维矢量记号表示每个理论的能量-动量张量的分量，使用散度、梯度、旋度、电场和磁场，并用上加点号表示时间导数。

\textbf{(b)}\quad 利用运动方程，验证（用你喜欢的任何记号）能量-动量张量是守恒的。

\textbf{13.}\quad 考虑在电磁学的拉格朗日量中添加一个形如 $\mathcal{L}' = \tilde{\epsilon}_{\mu\nu\rho\sigma}F^{\mu\nu}F^{\rho\sigma}$ 的附加项。

\textbf{(a)}\quad 用 $\mathbf{E}$ 和 $\mathbf{B}$ 表示 $\mathcal{L}'$。

\textbf{(b)}\quad 证明包含 $\mathcal{L}'$ 并不影响麦克斯韦方程组。你能想出一个深刻的理由吗？
